% !TeX root = ../../Book2.tex
% 纲要标题：## 第19章*　非交换局部化、多项式恒等式与增长
% 纲要位置：计划纲要.md，第 1668 行。

\OptionalChapter{非交换局部化、多项式恒等式与增长}{b2:ch:19}
\BookChapterNumber{b2:ch:19}

\begin{chapterabstract}
本章研究非交换环中允许分母换写的 Ore 条件，构造正则分母的局部化，并证明半素右 Goldie 环取得半单经典分式环的条件。矩阵的 Amitsur–Levitzki 多项式恒等式与有限生成代数的 Gelfand–Kirillov 增长维数，分别从恒等式和增长两方面考察非交换代数。
\end{chapterabstract}

\begin{learninggoals}
\begin{itemize}
\item 区分左右 Ore 方程，核对共同分母、分式等价关系、乘法和泛性质中的次序。
\item 掌握本质子模与一致维数，完整证明 Goldie 定理的充分方向，说明半素性及两项右 Goldie 条件的用途。
\item 证明矩阵的标准恒等式，说明有理系数证明怎样通过整数泛型矩阵推广到任意特征。
\item 计算有限生成代数的增长，证明 GK 维数不依赖生成空间，并区分增长、PI 性和分式存在性。
\end{itemize}
\end{learninggoals}

设 \(R=k\langle x,y\rangle\) 为域 \(k\) 上的自由结合代数。非零元素 \(x\) 与 \(y\) 没有非零的共同右倍数：\(xR\) 中各个字都以 \(x\) 开头，\(yR\) 中各个字都以 \(y\) 开头。若想把所有非零元素当成右分母，就必须能为这两个元素找到非零的共同右倍数，因而在此处已经失败。非交换分式不能仅靠“分母非零”来定义；分母能否越过分子，是必须先解决的乘法问题。

\ProseParagraphBreak

Ore 条件把分母换写要求表达为环内的方程；Goldie 理论则寻找使这些方程普遍可解的右理想条件，并判断分式环何时半单。本章另有两类相对独立的问题：一个非交换多项式能否在某个代数中对所有代入都取零，以及长度有限的乘积能张成多大的空间。前者检查所有代入共同满足的乘法关系，后者计算短字张成的线性空间。分母方程是否可解、是否存在这样的共同恒等式、短字怎样增长，是三个需要分别研究的问题；自由代数在分母方程处失败，并不把后两类问题变成前一类问题的推论。

\ProseParagraphBreak
非交换分式环与 Goldie 理论的进一步阅读，可参阅 T. Y. Lam 的《Lectures on Modules and Rings》\cite[\S 10A--10B, ``Ore Localizations'' and ``Right Ore Rings and Domains''; \S 11, ``Right Goldie Rings and Goldie's Theorems'']{Lam1999ModulesRings}。不同局部化理论对分母是否允许零因子、原环能否单射地映入局部化有不同要求。

% !TeX root = ../../Book2.tex
% 纲要标题：### 19.1 Ore 局部化
% 纲要位置：计划纲要.md，第 1670 行。

\section{Ore 局部化}
\label{b2:sec:1901}

交换环中把分母移到乘积的一侧，依靠的是元素可以交换。非交换环里，写下 \(as^{-1}\) 之后，还须解释 \(s^{-1}b\) 如何重新写成分母在右侧的形式。Ore 条件正是为这一步提供可解的方程。本节只讨论由双侧非零因子组成的分母集，从而原环能够单射到分式环。\cref{b2:prop:zero-divisor-inversion}将说明，若把原环中的零因子变成可逆元，单射性就必然丢失。

% 纲要要点（供目录核对）：
%
% * 分母在非交换情形的障碍；由实际右分式乘积产生 Ore 方程，正文证明只逆一个正则元素也会失败、逆零因子必破坏单射性
% * 左、右 Ore 条件；按最终分母位置区分，共同分母属于分母集而正则乘子不必属于其中
% * 典型局部化及其泛性质；先构造加法分式类，再由左乘算子复合给出结合乘法，解释被迫可逆的乘子及非中心矩阵计算
%

\subsection{分母在非交换情形的障碍}
\label{b2:subsec:1901:01}

本节设 \(R\ne0\) 为含单位元的环。若 \(s\in R\) 满足
\[
sx=0\Longrightarrow x=0,\qquad xs=0\Longrightarrow x=0,
\]
称 \(s\) 为\term[zhengze-yuansu]{正则元素}[regular element]，即它既不是左零因子，也不是右零因子。正则元素的乘积仍正则，一也是正则元素。

\ProseParagraphBreak
假设 \(R\) 已嵌入某个环 \(Q\)，分母集 \(S\) 中的元素在 \(Q\) 中可逆，并且每个 \(Q\) 元素都能写为右分式。对 \(p,a\in R\) 及 \(s,r\in S\)，乘积为
\[
(ps^{-1})(ar^{-1})=p(s^{-1}a)r^{-1}.
\]
要使中间的 \(s^{-1}a\) 也成为右分式，必须将它写成 \(bt^{-1}\)，其中 \(b\in R,t\in S\)。这个等式两侧分别左乘 \(s\)、右乘 \(t\)，就得到
\[
at=sb.
\]
因此共同倍数方程正是乘法中移换分母的要求。取得这样的 \(b,t\) 后，乘积便能整理为 \((pb)(rt)^{-1}\)。

\ProseParagraphBreak
以自由代数 \(R=k\langle x,y\rangle\) 为例。由\cref{b2:prop:free-associative-algebra}，有限字构成基。对非零多项式按字长、同长度内再按字典序选最大字，乘积中字长达到最大时，两因子都须各自达到最大字长；在这些固定长度的字中，任一因子按字典序变小，连接后的字也变小。因此乘积的最大字只能由两个最大字连接而成，系数为两个非零首系数的乘积，故乘积非零。因此所有非零元素正则。

\ProseParagraphBreak
但 \(xR\cap yR=0\)，因为第一项中的非零字都以 \(x\) 开头，第二项中的非零字都以 \(y\) 开头。若 \(a=x,s=y\) 能满足上述关系且 \(t\ne0\)，就有 \(0\ne xt=yb\in xR\cap yR\)，矛盾。所以这个非交换整环不能通过“所有非零元素作右分母，且每个元素是一个右分式”的方式得到分式除环。这不排除其他类型的扩张环，只说明通常的单分式表达需要额外条件。即使把分母限制为一个正则元素 \(x\) 的各次幂，所需方程也未必可解，见\cref{b2:prop:single-element-ore-obstruction}。

\subsection{左、右 Ore 条件}
\label{b2:subsec:1901:02}

\begin{definition}\label{b2:def:ore-set}
设 \(R\ne0\) 为含幺环，\(S\subseteq R\) 含一、乘法封闭，且全部元素均双侧正则。对任意 \(a\in R,s\in S\)，若能取 \(t\in S,b\in R\) 使 \(at=sb\)，则称 \(S\) 满足\term[you-Ore-tiaojian]{右 Ore 条件}[right Ore condition]，并称这样的 \(S\) 为\term[zhengze-you-Ore-ji]{正则右 Ore 集}[regular right Ore set]；在本节的正则情形，也称右 Ore 分母集。若对任意 \(a\in R,s\in S\) 都能取 \(t\in S,b\in R\) 使 \(ta=bs\)，则称 \(S\) 为正则左 Ore 集。两个方程及其分式方向为
\[
\begin{aligned}
\text{右 Ore：}&\quad at=sb,\qquad as^{-1};\\
\text{左 Ore：}&\quad ta=bs,\qquad s^{-1}a.
\end{aligned}
\]
这里的左右按最终分式中分母的位置命名。左条件就是反环中的右条件；即使分母都正则，两条件也不自动相同。
\end{definition}

\begin{lemma}\label{b2:lem:ore-common-denominators}
设 \(R\ne0\) 为含幺环，\(S\subseteq R\) 为正则右 Ore 集。
\begin{enumerate}
\item 若 \(s,su\in S\)，则 \(u\) 正则，但不一定属于 \(S\)。
\item 任意 \(s,t\in S\) 有共同右倍数 \(su=tv\in S\)，其中乘子 \(u,v\) 均正则。任意有限多个分母也有这样的共同右倍数。
\end{enumerate}
\end{lemma}
\begin{proof}
若 \(ux=0\)，则 \((su)x=0\)，由 \(su\) 正则得 \(x=0\)。对分子 \(s\) 和分母 \(su\) 使用 Ore 条件，取 \(v\in S,b\in R\) 使 \(sv=(su)b\)。消去左侧正则因子 \(s\)，得 \(v=ub\)。若 \(xu=0\)，则 \(xv=xub=0\)，由 \(v\) 正则得 \(x=0\)。故 \(u\) 双侧正则。

对分子 \(t\)、分母 \(s\) 使用 Ore 条件，得 \(tv=su\)，其中 \(v\in S\)。共同值属于 \(S\)，第一部分保证 \(u\) 也正则。有限多个分母逐次与已经取得的共同分母再作此操作即可；每个旧分母都仍右整除新的共同值。
\end{proof}

共同分母必须属于 \(S\)，放大分母所用的乘子则只需正则，以便在比较分子时消去。例如在 \(\mathbb Z\) 中取 \(S=\{2^i6^j:i,j\ge0\}\)，有 \(2,6\in S\) 及 \(2\cdot3=6\)，但乘子三不属于 \(S\)。

\ProseParagraphBreak
在交换环中，或更一般地在分母都位于中心时，取 \(t=s\)、\(b=a\) 即满足右 Ore 条件，左条件也成立。一般情形的共同倍数则可能需要非平凡的乘子；不能直接把共同分母写成 \(st\)，因为 \(t\) 未必右整除 \(st\)。

\subsection{典型局部化及其泛性质}
\label{b2:subsec:1901:03}

共同右分母足以定义分式的加法，而非交换乘法还要处理分母与分子相遇时的次序。先构造加法群，再将分式实现成这个群上的左乘算子，就能用算子复合取得结合乘法。若 \(su=tv=w\in S\)，预期的换分母等式是 \(as^{-1}=(au)w^{-1}\)、\(bt^{-1}=(bv)w^{-1}\)；因此把两分式视为相同，正是要求它们通分后有相同的分子。

\begin{lemma}[右分式类的加法]\label{b2:lem:ore-additive-classes}
设 \(R\ne0\) 为含幺环，\(S\subseteq R\) 为正则右 Ore 集。在 \(R\times S\) 上规定
\begin{equation}\label{b2:eq:ore-equivalence}
(a,s)\sim(b,t)
\quad\Longleftrightarrow\quad
\exists u,v\in R:\ su=tv\in S,\quad au=bv.
\end{equation}
则 \(\sim\) 是等价关系。其商集 \(F=(R\times S)/{\sim}\) 上的公式
\[
[a,s]+[b,t]=[au+bv,w]\quad(su=tv=w\in S),
\qquad -[a,s]=[-a,s]
\]
给出交换群结构。其零元满足
\begin{equation}\label{b2:eq:ore-zero-numerator}
[a,s]=0\quad\Longleftrightarrow\quad a=0.
\end{equation}
\end{lemma}
\begin{proof}
自反性与对称性直接成立。证明传递性时，设 \(su=tv=w\)、\(au=bv\)，又有 \(tu'=rv'=w'\)、\(bu'=cv'\)。取共同右倍数 \(w\alpha=w'\beta\in S\)。从 \(tv\alpha=tu'\beta\) 消去左侧正则元 \(t\)，得到 \(v\alpha=u'\beta\)，于是
\[
s u\alpha=r v'\beta,\qquad
a u\alpha=bv\alpha=bu'\beta=cv'\beta.
\]
故 \((a,s)\sim(c,r)\)。

先核对一个换分母事实。若 \((a,s)\sim(b,t)\)，则对任意共同右分母 \(w=su=tv\in S\)，都有 \(au=bv\)。事实上，取等价关系中原有的共同分母 \(w_0=su_0=tv_0\)，再取 \(w_0\alpha=w\beta\in S\)。消去 \(s,t\) 得 \(u_0\alpha=u\beta\)、\(v_0\alpha=v\beta\)，故
\[
(au-bv)\beta=(au_0-bv_0)\alpha=0.
\]
由\cref{b2:lem:ore-common-denominators}，乘子 \(\beta\) 正则，因而 \(au=bv\)。

固定两对代表元，若分别选 \(w=su=tv\) 与 \(w'=su'=tv'\)，取 \(W=w\alpha=w'\beta\in S\)。消去 \(s,t\) 得 \(u\alpha=u'\beta\) 及 \(v\alpha=v'\beta\)，所以两个候选和扩张到 \(W\) 后分子相同。这证明共同分母的选择无关。

若把第一对代表元 \((a,s)\) 换为等价的 \((a',s')\)，而 \((b,t)\) 保持不变，分别取 \(w=su=tv\)、\(w'=s'u'=tv'\)，再取 \(W=w\alpha=w'\beta\in S\)。由刚证的换分母事实，\(a u\alpha=a'u'\beta\)；由 \(tv\alpha=tv'\beta\) 消去 \(t\)，得 \(v\alpha=v'\beta\)。因此
\[
(au+bv)\alpha=(a'u'+bv')\beta,
\]
两个候选和等价。若改换第二对代表元 \((b,t)\sim(b',t')\)，固定 \((a,s)\)，分别取 \(w=su=tv\)、\(w'=su'=t'v'\)，再取 \(W=w\alpha=w'\beta\in S\)。消去 \(s\) 得 \(u\alpha=u'\beta\)；换分母事实给出 \(bv\alpha=b'v'\beta\)。于是
\[
(au+bv)\alpha=(au'+b'v')\beta,
\]
两个候选和仍等价。依次换两对代表元即可处理任意代表元选择。

由\cref{b2:lem:ore-common-denominators}，三个类可取同一共同右分母，结合律与交换律于是归结为 \(R\) 的加法；\([0,1]\) 是零元，\([-a,s]\) 是 \([a,s]\) 的负元。最后，\([a,s]=[0,1]\) 时存在 \(u,v\in R\) 使 \(su=v\in S\)、\(au=0\)；由 \(u\) 正则得 \(a=0\)。反向由共同分母 \(s\) 直接得到。
\end{proof}

\begin{theorem}[正则分母的右 Ore 局部化]\label{b2:thm:right-ore-localization}
设 \(R\ne0\) 为含幺环，\(S\subseteq R\) 为正则右 Ore 集，则存在含幺环 \(Q=RS^{-1}\) 及保幺单射 \(R\rightarrow Q\)，使 \(S\) 中每个元素在 \(Q\) 中可逆，且每个 \(Q\) 元素都可写为 \(as^{-1}\)，其中 \(a\in R,s\in S\)。对任意保幺环同态 \(f:R\rightarrow B\)，若全部 \(f(s)\) 在 \(B\) 中可逆，则存在唯一保幺环同态
\[
\widetilde f:Q\longrightarrow B,\qquad
\widetilde f(as^{-1})=f(a)f(s)^{-1}.
\]
\end{theorem}
\begin{proof}
由\cref{b2:lem:ore-additive-classes}得到右分式类的加法群 \(F\)，在它的自同态环 \(\operatorname{End}_{\mathbb Z}(F)\) 中构造分式环。

对 \(a\in R\) 定义加法群自同态 \(L_a[b,t]=[ab,t]\)。等价关系和加法公式说明它良定义且可加，且 \(L_{a+b}=L_a+L_b\)、\(L_aL_b=L_{ab}\)、\(L_1=\operatorname{id}\)。对分母 \(s\in S\)，还须证明 \(L_s\) 是自同构。它单射，因为 \([sb,t]=0\) 迫使 \(sb=0\)，进而 \(b=0\)。它也满射：对 \([b,t]\)，由 Ore 条件取 \(v\in S,u\in R\) 使 \(bv=su\)。因 \(tv\in S\)，可以通分为
\[
[b,t]=[bv,tv]=L_s[u,tv].
\]
故 \(L_s\) 可逆。

在 \(\operatorname{End}_{\mathbb Z}(F)\) 中考虑所有算子 \(L_aL_s^{-1}\)。若 \(su=tv=w\in S\)，由 \(L_w=L_sL_u=L_tL_v\) 可知 \(L_u=L_s^{-1}L_w\)、\(L_v=L_t^{-1}L_w\) 也可逆。这一步用的是 \(s,t,w\) 都已被逆转，而不只是乘子 \(u,v\) 正则。于是
\[
L_aL_s^{-1}=L_{au}L_w^{-1},\qquad
L_bL_t^{-1}=L_{bv}L_w^{-1}.
\]
因此它们对加法与负元封闭，且等价分式给出相同算子。若 \(bv=su\)、\(v\in S\)，则
\[
L_s^{-1}L_b=L_uL_v^{-1},
\]
从而
\begin{equation}\label{b2:eq:right-ore-product}
\begin{aligned}
(L_aL_s^{-1})(L_bL_t^{-1})
&=L_{au}L_v^{-1}L_t^{-1}\\
&=L_{au}L_{tv}^{-1}.
\end{aligned}
\end{equation}
所以这些算子组成含单位元的子环。它们与 \(F\) 的元素一一对应：\(L_aL_s^{-1}\) 在 \([1,1]\) 处的值恰为 \([a,s]\)，因为 \(L_s[1,s]=[1,1]\)。故相同算子必来自相同分式类，反向已由共同分母证明。将这个子环的乘法转移到 \(F\)，得到结合环 \(Q\)。乘法结合律由算子复合继承，分配律也不需另作分母猜测。

映射 \(a\mapsto L_a\) 为保单位元环同态，且由在 \([1,1]\) 处取值及式\eqref{b2:eq:ore-zero-numerator}可知它单射。每个 \(L_s^{-1}=L_1L_s^{-1}\) 也在 \(Q\) 中，所以分母可逆。于是将 \([a,s]\) 记作 \(as^{-1}\) 后，式\eqref{b2:eq:right-ore-product}给出具体乘法
\[
(as^{-1})(bt^{-1})=(au)(tv)^{-1},
\qquad bv=su,\quad v\in S.
\]

最后核对泛性质。若 \(su=tv=w\in S\)，则 \(f(u)=f(s)^{-1}f(w)\) 可逆：即使 \(u\notin S\)，任何把 \(s,w\) 变成单位的同态也会迫使 \(u\) 变成单位。并有
\[
f(a)f(s)^{-1}=f(au)f(w)^{-1}.
\]
因此式\eqref{b2:eq:ore-equivalence}的相同分子保证 \(\widetilde f\) 良定义，共同分母公式保证可加。由 \(bv=su\)，得到 \(f(s)^{-1}f(b)=f(u)f(v)^{-1}\)，恰好验证乘法公式被保持。单位也被保持。每个元素都是 \(as^{-1}\)，其像已由 \(f\) 及分母求逆决定，故延拓唯一。
\end{proof}
\symidx{\(RS^{-1}\)：右 Ore 分式环}

这称为\term[Ore-jubuhua]{Ore 局部化}[Ore localization]。对反环应用同一定理，得到左 Ore 分式环 \(S^{-1}R\)。若同一个 \(S\) 同时满足左右条件，两种环都满足使 \(S\) 中元素可逆的同一个泛性质。分别延拓原环到另一分式环的映射，得到两个方向的同态；其复合都在 \(R\) 上为恒等，泛性质中的唯一性便使复合在整个分式环上也为恒等。因此两种构造具有保持 \(R\) 的互逆同构。这里所用的唯一性论证与\cref{b2:prop:universal-object-uniqueness}中积、余积的情形相同。

\begin{proposition}\label{b2:prop:ore-clearing-denominators}
设 \(R\ne0\) 为含幺环，\(S\subseteq R\) 为正则右 Ore 集，\(Q=RS^{-1}\) 为其右 Ore 局部化。对任意有限组 \(q_1,\ldots,q_m\in Q\)，存在共同右分母 \(s\in S\) 使全部 \(q_i s\in R\)。若 \(I\) 为 \(R\) 的右理想，则
\[
IQ=\{\,as^{-1}:a\in I,\ s\in S\,\}.
\]
每个非零右 \(R\) 子模 \(X\subseteq Q\) 都与 \(R\) 有非零交。
\end{proposition}
\begin{proof}
先把各分母取共同右倍数，分子随之扩张，就得到第一项。\(IQ\) 中元素为有限和 \(\sum_i a_iq_i\)、\(a_i\in I\)，通分后分子仍属于右理想 \(I\)，得到第二式的一向；另一向由 \(as^{-1}\in IQ\) 直接成立。最后，若 \(0\ne as^{-1}\in X\)，右乘 \(s\in R\) 得到 \(0\ne a\in X\cap R\)，其中非零性来自分母在 \(Q\) 中可逆。
\end{proof}

例如令 \(C\) 为交换整环，\(R=M_n(C)\)，分母取非零标量矩阵 \(sI_n\)。它们在中心且正则，局部化得到
\[
RS^{-1}\cong M_n(\operatorname{Frac}C).
\]
\symidx{\(\operatorname{Frac}C\)：交换整环的分式域}
每个右分式逐条目取标量分母，反向对一个分式域矩阵的有限多个条目取共同分母，便写成一个矩阵右分式；两种操作互逆且保持加乘。原环可能有许多矩阵零因子，但这些元素没有被放入分母集。

\ProseParagraphBreak
也可以让分母不在中心，此时 \(s^{-1}a\) 与 \(as^{-1}\) 可能不同，换分母还须改变分子。取 \(R=M_2(\mathbb Z)\)，令 \(S\) 为其中行列式非零的全部矩阵。它们在 \(M_2(\mathbb Q)\) 中可逆，因此在 \(R\) 中双侧正则，且 \(S\) 含一、乘法封闭。给定 \(a\in R,s\in S\)，选正整数 \(q\) 清除 \(s^{-1}a\) 的全部条目分母，令 \(t=qI_2\in S\)、\(b=s^{-1}at\in R\)，便有 \(at=sb\)。所以 \(S\) 满足右 Ore 条件。所得分式环仍为 \(M_2(\mathbb Q)\)：每个右分式可在该矩阵环中解释，而每个有理矩阵都能用非零整数标量矩阵作共同分母；若右分式在其中为零，其分子也必为零。

\ProseParagraphBreak
具体取
\[
s=\begin{pmatrix}1&1\\0&2\end{pmatrix},\qquad
a=E_{21},\qquad t=2I_2,\qquad
b=\begin{pmatrix}-1&0\\1&0\end{pmatrix}.
\]
直接相乘得到
\[
at=\begin{pmatrix}0&0\\2&0\end{pmatrix}=sb.
\]
于是非中心分母可以按 Ore 方程移到右侧，但分子也随之改变：
\[
s^{-1}a=bt^{-1}
=\begin{pmatrix}-\frac12&0\\\frac12&0\end{pmatrix},
\qquad
as^{-1}=\begin{pmatrix}0&0\\1&-\frac12\end{pmatrix}.
\]
例如两个右分式的乘积为
\[
\begin{aligned}
(E_{11}s^{-1})(as^{-1})
&=E_{11}b\,t^{-1}s^{-1}\\
&=E_{11}b(st)^{-1}\\
&=-E_{11}(2s)^{-1}
=\begin{pmatrix}-\frac12&\frac14\\0&0\end{pmatrix}.
\end{aligned}
\]
这里先用 \(s^{-1}a=bt^{-1}\) 处理两个分式之间相遇的分母与分子，再用 \(t^{-1}s^{-1}=(st)^{-1}\) 合并分母；次序来自实际乘法。

\begin{proposition}\label{b2:prop:single-element-ore-obstruction}
设 \(k\) 为域，\(R=k\langle x,y\rangle\) 为两个生成元上的自由结合含幺 \(k\) 代数，\(S=\{x^n:n\ge0\}\)。则 \(S\) 含一、乘法封闭，全部元素双侧正则，但不满足右 Ore 条件。
\end{proposition}
\begin{proof}
由\cref{b2:prop:free-associative-algebra}，\(x,y\) 上的有限字是 \(R\) 的一组 \(k\) 基，空字表示一。对每个 \(n\ge0\)，在字的左侧或右侧连接 \(x^n\)，都把不同的字送到不同的字。因此左乘与右乘 \(x^n\) 都是单射，\(x^n\) 双侧正则。\(x^0=1\) 及 \(x^i x^j=x^{i+j}\) 还说明 \(S\) 含一、乘法封闭。

若满足右 Ore 条件，对 \(a=y,s=x\) 就能取 \(t=x^n\in S\)、\(b\in R\)，使 \(yx^n=xb\)。左侧是以 \(y\) 开头的非零基字，右侧的每个非零基字都以 \(x\) 开头。字基的线性无关性排除该等式，所以右 Ore 条件失败。
\end{proof}

\begin{proposition}\label{b2:prop:zero-divisor-inversion}
设 \(R\ne0\)、\(B\) 为含幺环，\(f:R\to B\) 为保幺单射。若 \(s\in R\) 的像 \(f(s)\) 可逆，则 \(s\) 双侧正则。

具体地，设 \(k\) 为域，\(R=k\times k\)、\(e=(1,0)\)。第一坐标投影 \(\pi:R\to k\) 把 \(e\) 送到单位元，且对每个含幺环 \(B\) 及每个使 \(f(e)\) 可逆的保幺环同态 \(f:R\to B\)，都存在唯一保幺环同态 \(\bar f:k\to B\) 使 \(f=\bar f\pi\)。但是 \(\ker\pi=0\times k\ne0\)，任何这样的 \(f\) 都不是单射。
\end{proposition}
\begin{proof}
若 \(sx=0\)，则 \(f(s)f(x)=0\)，乘以 \(f(s)^{-1}\) 得 \(f(x)=0\)，再由单射性得 \(x=0\)。若 \(xs=0\)，在右侧乘以 \(f(s)^{-1}\) 同样得到 \(x=0\)。所以可嵌入后变成单位的元素必须双侧正则。

在第二个例子中，\(e^2=e\)。可逆的 \(f(e)\) 满足 \(f(e)^2=f(e)\)，消去一个因子得到 \(f(e)=1_B\)，于是 \(f(1-e)=0\)。每个第二坐标元素都为 \((0,b)=(1-e)(0,b)\)，故被 \(f\) 送到零。定义 \(\bar f(a)=f(a,0)\)，它保持加法与乘法，并把一送到 \(f(e)=1_B\)，因而为保幺环同态。又有 \(f(a,b)=f(a,0)=\bar f(\pi(a,b))\)，所以 \(f=\bar f\pi\)。投影 \(\pi\) 满射，保证这种分解唯一；它本身满足 \(\pi(e)=1\)，且核恰为非零的第二坐标因子。因此使 \(e\) 可逆必然零化原环中的非零元素。
\end{proof}

% !TeX root = ../../Book2.tex
% 纲要标题：### 19.2 Goldie 理论
% 纲要位置：计划纲要.md，第 1676 行。

\section{Goldie 理论}
\label{b2:sec:1902}

Ore 条件逐个要求方程 \(at=sb\) 可解，直接核验往往困难。Goldie 理论用环的右理想结构保证这些方程有解，并进一步说明所得分式环为何半单。本节明确改用右模：一致维数、子模和零化子链均按所写的作用方向理解。证明先建立本质右理想中正则元素的存在性，再进行局部化。

% 纲要要点（供目录核对）：
%
% * 一致维数及零化子条件；以整数、向量空间和截断多项式模区别本质、一致、长度，解释有限本质一致直和的构造与逆向证明
% * 半素右 Goldie 环的分式环；分别追踪半素性、右零化子升链条件及有限一致维数，补清本质理想含正则元、Ore 乘子的原像和右半单转为左半单的步骤
% * 定理假设与交换整环的比较；正文证明整环右 Ore 等价于右正则模一致，正式证明上三角环中的本质理想无正则元，区分矩阵的行模数与基域维数
%

\subsection{一致维数与零化子条件}
\label{b2:subsec:1902:01}

\begin{definition}\label{b2:def:essential-uniform-dimension}
设 \(R\) 为含幺环，\(M\) 为幺性右 \(R\) 模，\(E\subseteq M\) 为子模。若每个非零子模 \(N\subseteq M\) 都满足 \(N\cap E\ne0\)，则称 \(E\) 为 \(M\) 的\term[benzhi-zimo]{本质子模}[essential submodule]。若非零幺性右 \(R\) 模 \(U\) 的任意两个非零子模交非零，则称 \(U\) 为\term[yizhi-mo]{一致模}[uniform module]；等价地，\(U\) 的每个非零子模都在 \(U\) 中本质。

若 \(M\) 的非零子模 \(N_1,\ldots,N_r\) 的和为内部直和，则称它们独立。定义 \(M\) 的\term[yizhi-weishu]{一致维数}[uniform dimension] \(\operatorname{udim}M\) 为这类有限独立族大小的上确界，允许为无穷；零模的值为零。一致维数有限时也称为 \term[Goldie-zhi]{Goldie 秩}[Goldie rank]。
\end{definition}
\symidx{\(\operatorname{udim}M\)：右模的一致维数}

本质包含要求每个非零子模都与给定子模相交，并不要求给定子模等于全模。例如 \(2\mathbb Z\) 在 \(\mathbb Z\) 中本质，因为任意非零子模 \(d\mathbb Z\) 都与它有非零公共元素 \(2d\)。但域 \(k\) 上向量空间 \(V\) 的真子空间 \(W\) 不本质：取 \(v\notin W\)，则 \(kv\cap W=0\)。

\ProseParagraphBreak
一致模不能同时容纳两个交零的非零子模，所以对非零模 \(U\)，一致性等价于 \(\operatorname{udim}U=1\)。对除环上的右向量空间，这恰好是维数为一的条件：一维空间只有零子空间与自身，而两个线性无关向量张成的直线交于零。简单模因此一致，但一致模未必简单。例如交换整环 \(R\) 的右正则模一致：两个非零理想分别取非零元素 \(a,b\)，非零乘积 \(ab\) 落在两者的交中。整数模因而一致，虽然有无限严格降链 \(\mathbb Z\supsetneq2\mathbb Z\supsetneq4\mathbb Z\supsetneq\cdots\)，并不具有有限长度。

\ProseParagraphBreak
一致维数数的是同时独立的直和项，合成长度则还计算这些项内部的简单层。取域 \(k\)、整数 \(m\ge1\) 及 \(R=k[t]/(t^m)\)，把 \(t\) 的剩余类仍记为 \(t\)。右理想与 \(k[t]\) 中包含 \((t^m)\) 的理想对应；\(k[t]\) 为主理想整环，这些理想的生成元整除 \(t^m\)，故可取为 \(t^i\)。所以 \(R\) 的右理想恰为 \(t^iR\)、\(0\le i\le m\)。它们彼此可比，任意两个非零右理想的交非零，因此 \(\operatorname{udim}R_R=1\)。另一方面，链
\[
0=t^mR\subsetneq t^{m-1}R\subsetneq\cdots\subsetneq tR\subsetneq R
\]
的每个相邻商 \(t^iR/t^{i+1}R\) 都由 \(t^i\) 的类张成，且 \(t\) 在该商上作用为零，所以同构于简单右模 \(R/tR\cong k\)。这给出长度为 \(m\) 的合成列，故 \(\ell(R_R)=m\)。

\begin{lemma}\label{b2:lem:essential-tools}
设 \(R\) 为含幺环，所涉右 \(R\) 模均为幺性模。若 \(M\) 的子模 \(E\subseteq N\subseteq M\) 满足 \(E\) 在 \(N\) 中本质、\(N\) 在 \(M\) 中本质，则 \(E\) 在 \(M\) 中本质。若 \(E\subseteq N\) 是右 \(R\) 模中的本质包含，\(f:M\rightarrow N\) 是右 \(R\) 模同态，则 \(f^{-1}(E)\) 在 \(M\) 中本质。对有限族右 \(R\) 模中的本质包含 \(E_i\subseteq M_i\)，还有本质包含 \(\bigoplus_iE_i\subseteq\bigoplus_iM_i\)。
\end{lemma}
\begin{proof}
若 \(E\) 本质于 \(N\)，而 \(N\) 本质于 \(M\)，任意非零 \(L\subseteq M\) 先与 \(N\) 非零相交，再与 \(E\) 非零相交，故传递性成立。

对非零子模 \(L\subseteq M\)，若 \(f(L)=0\)，则 \(L\subseteq f^{-1}(E)\)；否则 \(f(L)\cap E\ne0\)，取一个非零像的原像 \(x\in L\)，便有 \(0\ne x\in L\cap f^{-1}(E)\)。这证明原像的本质性。

在有限直和中，第 \(i\) 个投影的原像 \(\pi_i^{-1}(E_i)\) 由上段本质。有限个本质子模的交仍本质：先与第一个相交，再逐个与其余相交，每一步保持非零。它们的交恰为 \(\bigoplus_iE_i\)，得到最后结论。
\end{proof}

\begin{lemma}\label{b2:lem:uniform-submodule-existence}
设 \(R\) 为含幺环，\(M\) 为非零幺性右 \(R\) 模。若 \(M\) 的一致维数有限，或 \(M\) 为 Noether 模，则 \(M\) 的每个非零子模都含有一致子模。
\end{lemma}
\begin{proof}
假设某个非零子模 \(N\subseteq M\) 不含一致子模。令 \(N_0=N\)。由于 \(N_0\) 不一致，可取非零子模 \(N_1,W_0\subseteq N_0\)，使 \(N_1\cap W_0=0\)。\(N_1\) 也不含一致子模，故可继续在 \(N_1\) 中取非零子模 \(N_2,W_1\)，使 \(N_2\cap W_1=0\)，依此类推。每一步都有 \(N_{i+1}\oplus W_i\subseteq N_i\)，归纳可知 \(W_0,W_1,\ldots\) 独立。

若 \(\operatorname{udim}M<\infty\)，任意多个 \(W_i\) 构成的有限独立族便与一致维数的上界矛盾。若 \(M\) 为 Noether 模，则 \(W_0\subsetneq W_0\oplus W_1\subsetneq\cdots\) 是子模的无限严格升链，同样矛盾。
\end{proof}

\begin{proposition}\label{b2:prop:uniform-dimension-structure}
设 \(R\) 为含幺环，\(M\) 为非零幺性右 \(R\) 模。若 \(\operatorname{udim}M=d<\infty\)，则 \(M\) 含有本质子模 \(E=U_1\oplus\cdots\oplus U_d\)，其中各 \(U_i\) 是一致右 \(R\) 模。反过来，只要 \(M\) 含有这样的有限本质直和，就有 \(\operatorname{udim}M=d\)。若子模 \(N\subseteq M\) 且 \(M\) 一致维数有限，则
\[
\operatorname{udim}N=\operatorname{udim}M
\quad\Longleftrightarrow\quad N\text{ 在 }M\text{ 中本质}.
\]
每个 Noether 右 \(R\) 模都有有限一致维数。
\end{proposition}
\begin{proof}
先设 \(\operatorname{udim}M=d<\infty\)。由\cref{b2:lem:uniform-submodule-existence}取一致子模 \(U_1\)。若已取的有限直和 \(U_1\oplus\cdots\oplus U_r\) 不本质，就有非零子模与它交于零，在其中再取一致子模 \(U_{r+1}\)。一致维数限制新增项数，所以过程在有限步后停止，得到本质直和 \(E\)。

现在证明：本质直和 \(E=\bigoplus_{i=1}^dU_i\) 存在时，任意独立非零族 \(N_1,\ldots,N_s\subseteq M\) 都满足 \(s\le d\)。先以 \(N_j\cap E\) 代替 \(N_j\)，这些交仍非零且独立，故可假设它们在 \(E\) 中。对 \(d\) 归纳，\(d=0\) 时没有非零子模。若 \(s>0\)，令 \(W=N_2\oplus\cdots\oplus N_s\)。它与非零 \(N_1\) 交于零，所以不在 \(E\) 中本质。要减少一致直和的项数，可以找一个与 \(W\) 交零的 \(U_i\)，再将它商去；关键是证明这样的 \(U_i\) 必存在。如果每个 \(U_i\cap W\) 都非零，一致性及\cref{b2:lem:essential-tools}将使这些交的直和在 \(E\) 中本质，而该直和包含于 \(W\)，迫使 \(W\) 与 \(N_1\) 非零相交，矛盾。因此某个 \(U_i\cap W=0\)。商映射 \(E\rightarrow E/U_i\) 在 \(W\) 上单射，故其余 \(s-1\) 项在 \(E/U_i\cong\bigoplus_{j\ne i}U_j\) 中仍独立。归纳给出 \(s-1\le d-1\)。原来的 \(U_i\) 又给出大小 \(d\) 的独立族，故一致维数恰为 \(d\)。因此，前段构造的本质直和恰有 \(d\) 项。

若 \(M\) 为 Noether 模，\cref{b2:lem:uniform-submodule-existence}同样保证每次都能取出一致子模。若抽取过程不停，所得部分直和就构成严格升链 \(U_1\subsetneq U_1\oplus U_2\subsetneq\cdots\)，违反第十二章的子模升链条件。因此过程在有限步后停止，所得本质直和有有限项数；由刚证的逆向结论，\(M\) 的一致维数有限。

最后，若 \(N\) 本质，把 \(M\) 中任意最大独立族的每项与 \(N\) 相交，得到 \(N\) 中同样大的独立族，所以两维数相等。若 \(N\) 不本质，选 \(0\ne W\subseteq M\) 与 \(N\) 交零。将它加入 \(N\) 的一个最大独立族，得到 \(\operatorname{udim}M\ge\operatorname{udim}N+1\)，故两者不等。
\end{proof}

对 \(X\subseteq R\)，分别记
\[
\begin{aligned}
\operatorname{ann}_r(X)&=\{\,a\in R:xa=0\text{ 对所有 }x\in X\,\},\\
\operatorname{ann}_{\ell}(X)&=\{\,a\in R:ax=0\text{ 对所有 }x\in X\,\}.
\end{aligned}
\]
它们分别为右理想与左理想，称为右零化子与左零化子。
\symidx{\(\operatorname{ann}_r(X)\)：子集的右零化子}
\symidx{\(\operatorname{ann}_{\ell}(X)\)：子集的左零化子}

\begin{definition}\label{b2:def:semiprime-right-goldie}
设 \(R\) 为含幺环。若 \(R\) 没有非零的平方为零的双边理想，则称 \(R\) 为\term[bansu-huan]{半素环}[semiprime ring]。对 \(X\subseteq R\)，记 \(\operatorname{ann}_r(X)=\{\,r\in R\mid xr=0\text{ 对所有 }x\in X\,\}\)。若右正则模 \(R_R\) 一致维数有限，且这些右零化子组成的每条升链都最终稳定，则称 \(R\) 为\term[you-Goldie-huan]{右 Goldie 环}[right Goldie ring]。
\end{definition}

半素也等价于没有非零幂零双边理想：若 \(I^m=0\) 且 \(m\ge2\) 最小，则 \(I^{m-1}\ne0\) 而它的平方为零。这里排除的是双边理想，不是幂零元素；例如对域 \(k\)，单环 \(M_2(k)\) 半素，却有 \(E_{12}\ne0\)、\(E_{12}^2=0\)。半素环中的非零左理想或右理想，其平方也不能为零。比如左理想 \(L\) 若满足 \(L^2=0\)，双边理想 \(LR\) 满足
\((LR)^2\subseteq L^2R=0\)，故 \(LR=0\)，进而 \(L=0\)；右理想采用 \(RL\) 的对应论证。这里零化子升链条件只针对这些特殊右理想，比全部右理想都满足升链条件弱。

\subsection{半素右 Goldie 环的分式环}
\label{b2:subsec:1902:02}

下面要从半素右 Goldie 条件推出每个本质右理想都包含正则元素。半素性使非零一侧理想中能找到乘积非零的元素，右零化子的升链条件允许选择不能继续扩大的零化子；有限一致维数则限制可以加入的独立子模数目，使寻找单射左乘作用的过程在有限步后完成。三个条件将分别用于这些具体论证。

\begin{lemma}\label{b2:lem:essential-annihilator-zero}
设 \(R\) 为非零含幺半素右 Goldie 环，\(E\subseteq R_R\) 为本质右理想，则其左零化子 \(\operatorname{ann}_{\ell}(E)=\{\,r\in R\mid re=0\text{ 对所有 }e\in E\,\}\) 为零。
\end{lemma}
\begin{proof}
若存在本质右理想 \(E\) 满足 \(\operatorname{ann}_{\ell}(E)\ne0\)，取非零左理想 \(L=\operatorname{ann}_{\ell}(E)\)，则 \(\operatorname{ann}_r(L)\supseteq E\)，因而也是本质右理想。故下列右零化子的集合非空：
\[
\mathcal A=
\bigl\{\operatorname{ann}_r(L):
0\ne L\subseteq R\text{ 为左理想，且 }\operatorname{ann}_r(L)\text{ 在 }R_R\text{ 中本质}\bigr\}.
\]
沿用\cref{b2:prop:chain-maximal-minimal}证明中的升链论证，若 \(\mathcal A\) 没有极大元，就可在其中逐次严格增大，得到一条永不稳定的右零化子升链，违反右 Goldie 条件。因此可取 \(A=\operatorname{ann}_r(L)\in\mathcal A\) 为按包含关系极大的元素。

半素性给出 \(L^2\ne0\)，取 \(x,y\in L\) 使 \(xy\ne0\)。\(yR\) 是非零右理想，本质性使其中有非零元素 \(ya\in A\)。因为 \(x\in L\)、\(ya\in\operatorname{ann}_r(L)\)，有 \(xya=0\)，而 \(ya\ne0\)；也就是说，右乘子 \(a\) 零化了 \(xy\)，却没有零化 \(y\)。所以
\[
A\subseteq\operatorname{ann}_r(y)\subsetneq\operatorname{ann}_r(xy).
\]
这里 \(y,xy\in L\)，所以两者都包含 \(A\)；并且
\[
\operatorname{ann}_r(y)=\operatorname{ann}_r(Ry),\qquad
\operatorname{ann}_r(xy)=\operatorname{ann}_r(Rxy).
\]
由于 \(Ry,Rxy\) 均为非零左理想，而包含 \(A\) 的右理想仍本质，这两个零化子都属于 \(\mathcal A\)。\(A\) 的极大性要求它们都等于 \(A\)，与严格包含矛盾。
\end{proof}

\begin{lemma}\label{b2:lem:one-sided-regular-goldie}
设 \(R\) 为非零含幺半素右 Goldie 环，\(s\in R\)。若 \(\operatorname{ann}_r(s)=\{\,r\in R\mid sr=0\,\}=0\)，则 \(s\) 双侧正则，且 \(sR\) 为 \(R_R\) 的本质右理想。
\end{lemma}
\begin{proof}
左乘 \(s\) 给出右模同构 \(R_R\cong sR\)，所以两者一致维数相同。由\cref{b2:prop:uniform-dimension-structure}，\(sR\) 在 \(R\) 中本质。上一引理给出 \(\operatorname{ann}_{\ell}(sR)=0\)，而它正是 \(\operatorname{ann}_{\ell}(s)\)。两侧零化子均为零，所以 \(s\) 正则。
\end{proof}

\begin{lemma}\label{b2:lem:goldie-ideal-regular-element}
设 \(R\) 为非零含幺半素右 Goldie 环，则每个右理想 \(N\subseteq R\) 都含有某个 \(v\in N\)，使 \(\operatorname{ann}_r(v)\cap N=0\)。因此每个本质右理想都含有双侧正则元素。
\end{lemma}
\begin{proof}
先设 \(U\) 为非零一致右理想。由半素性，取 \(x,y\in U\) 使 \(xy\ne0\)。若 \(W=\operatorname{ann}_r(x)\cap U\ne0\)，则 \(W\) 在 \(U\) 中本质。对所有满足 \(ya\in W\) 的右乘子 \(a\)，都有 \(xya=0\)；为了把这些乘子组成的右理想与本质性联系起来，考虑右线性映射 \(\lambda_y:R\rightarrow U\)、\(a\mapsto ya\)。这些乘子恰好组成 \(\lambda_y^{-1}(W)\)，由\cref{b2:lem:essential-tools}它是本质右理想，却被非零元素 \(xy\) 从左零化，违反\cref{b2:lem:essential-annihilator-zero}。故 \(W=0\)，一致右理想中的结论成立。

若 \(N\ne0\)，它含有一致右理想，上一段因而保证存在满足
\[
v\in V\subseteq N,\qquad \operatorname{ann}_r(v)\cap V=0
\]
的非零子模 \(V\) 及其中的元素 \(v\)。选择一对 \((V,v)\)，使 \(\operatorname{udim}V\) 最大；这些维数是被 \(\operatorname{udim}R_R\) 限制的正整数，所以最大值存在。假设仍有 \(\operatorname{ann}_r(v)\cap N\ne0\)，在其中取一致右理想 \(U\)，再由上段取 \(u\in U\) 使 \(\operatorname{ann}_r(u)\cap U=0\)。因为 \(U\subseteq\operatorname{ann}_r(v)\)，有 \(U\cap V=0\)。

令 \(x=u+v\)，是为了在保留 \(v\) 对 \(V\) 的单射作用时，用 \(u\) 补上它在 \(U\) 上的零作用。这里 \(U,V\) 都是右理想，所以 \(u(U\oplus V)\subseteq U\)、\(vV\subseteq V\)，又有 \(vU=0\)。因此左乘 \(x\) 确实给出右模自同态 \(\lambda_x:U\oplus V\rightarrow U\oplus V\)，相对于这个分解写为
\[
\lambda_x=
\begin{pmatrix}
\lambda_u|_U&\lambda_u|_V\\
0&\lambda_v|_V
\end{pmatrix}.
\]
两条对角映射按 \(u,v\) 的选择都是单射，左下角则为零，所以相加不会产生新的核。具体地，若 \(z=a+b\)、\(a\in U,b\in V\)，且 \(xz=0\)，先看 \(V\) 分量得到 \(vb=0\)，从而 \(b=0\)；再看 \(U\) 分量得到 \(ua=0\)，从而 \(a=0\)。故 \(\operatorname{ann}_r(x)\cap(U\oplus V)=0\)。但 \(x\in U\oplus V\subseteq N\)，而把 \(V\) 中大小为 \(\operatorname{udim}V\) 的独立族与非零项 \(U\) 并在一起，就使一致维数至少增加一，与选择的最大值矛盾。因此 \(\operatorname{ann}_r(v)\cap N=0\)。\(N=0\) 时取 \(v=0\) 即可。

若 \(N\) 本质，这个零交迫使右理想 \(\operatorname{ann}_r(v)\) 为零。由\cref{b2:lem:one-sided-regular-goldie}，\(v\) 双侧正则。
\end{proof}

\begin{theorem}[Goldie 定理的充分方向]\label{b2:thm:goldie-sufficiency}
设 \(R\) 为非零含幺半素右 Goldie 环，\(S\subseteq R\) 为全部双侧正则元素，则 \(S\) 满足右 Ore 条件，经典右分式环 \(Q=RS^{-1}\) 为半单 Artin 环。
\end{theorem}
\begin{proof}
一属于 \(S\)，正则元素乘法封闭。给定 \(a\in R,s\in S\)，要解 \(at=sb\)，就是要找到正则的右乘子 \(t\)，使 \(at\in sR\)。这些候选乘子组成右模同态 \(\lambda_a:R_R\rightarrow R_R\)、\(x\mapsto ax\) 的原像 \(\lambda_a^{-1}(sR)\)。由\cref{b2:lem:one-sided-regular-goldie}，\(sR\) 本质，所以\cref{b2:lem:essential-tools}使这个原像也本质。再由\cref{b2:lem:goldie-ideal-regular-element}取其中的正则元素 \(t\)，便有 \(at=sb\) 对某个 \(b\in R\) 成立，得到右 Ore 条件。\cref{b2:thm:right-ore-localization}给出 \(Q\)。

局部化把正则元素变成单位，这将排除 \(Q\) 中的真本质右理想。设 \(E\) 为 \(Q\) 的本质右理想，\(I\) 为 \(R\) 的任意非零右理想，则 \(IQ\) 是非零右 \(Q\) 理想，所以 \(E\cap IQ\) 含有非零元素 \(q\)。由\cref{b2:prop:ore-clearing-denominators}，可取 \(s\in S\) 使 \(qs\in I\)，且 \(s\) 在 \(Q\) 中可逆保证 \(qs\ne0\)。因为 \(E\) 是右理想，\(qs\) 也在 \(E\)，因此 \(E\cap R\) 是 \(R\) 的本质右理想。它含有正则元素 \(t\)，而 \(t^{-1}\in Q\)；右理想 \(E\) 因而含有 \(tt^{-1}=1\)，只能等于 \(Q\)。

要将本质理想的结论变成直和分解，还需为任意右理想找到与它交零、且二者之和本质的子模。任意独立的非零右 \(Q\) 理想与 \(R\) 的交均非零，且作为右 \(R\) 子模仍独立，所以 \(\operatorname{udim}Q_Q\le\operatorname{udim}R_R<\infty\)。对给定右 \(Q\) 理想 \(J\)，在与 \(J\) 交零的右理想中选一致维数最大的 \(K\)：零理想保证候选非空，而这些维数是有上界的非负整数，保证最大值存在。若 \(J\oplus K\) 不本质，取非零右理想 \(L\) 与之交零；则 \(K\oplus L\) 仍与 \(J\) 交零，且将 \(L\) 加入 \(K\) 的最大独立族使一致维数至少增加一，违反所选最大值。故 \(J\oplus K\) 本质。刚证 \(Q\) 没有真本质右理想，所以 \(J\oplus K=Q\)。

因此右正则模 \(Q_Q\) 的每个子模都有补模。此前的\cref{b2:thm:semisimple-module-characterizations}以左模陈述，而半单环也按左正则模定义；把右 \(Q\) 模视为左 \(Q^{\mathrm{op}}\) 模，便可对同一个空间及其子模应用这一刻画，得到
\[
Q_Q\text{ 半单}
\quad\Longleftrightarrow\quad
{}_{Q^{\mathrm{op}}}Q\text{ 半单}
\quad\Longleftrightarrow\quad
Q^{\mathrm{op}}\text{ 为半单环}.
\]
对 \(Q^{\mathrm{op}}\) 应用\cref{b2:thm:artin-wedderburn}，存在有限多个除环 \(D_i\) 和正整数 \(n_i\)，使
\[
Q^{\mathrm{op}}\cong\prod_i M_{n_i}(D_i).
\]
取反环，并以矩阵转置识别 \(M_{n_i}(D_i)^{\mathrm{op}}\cong M_{n_i}(D_i^{\mathrm{op}})\)，便得
\[
Q\cong\prod_i M_{n_i}(D_i^{\mathrm{op}}).
\]
故 \(Q\) 是半单 Artin 环。
\end{proof}

这给出本章所用的\term[Goldie-dingli]{Goldie 定理}[Goldie's theorem]：由半素右 Goldie 条件得到半单经典分式环。本章证明这一充分方向，反向刻画留待进一步阅读。

\begin{corollary}\label{b2:cor:goldie-rank}
设 \(R\) 为非零含幺半素右 Goldie 环，\(S\subseteq R\) 为全部双侧正则元素，\(Q=RS^{-1}\) 为经典右分式环，则 \(\operatorname{udim}R_R=\ell(Q_Q)\)。
\end{corollary}
\begin{proof}
若 \(I_1,\ldots,I_d\) 为 \(R\) 中独立非零右理想，则 \(I_iQ\) 仍独立：一条和为零的关系中的各项，通到共同右分母后成为各 \(I_i\) 中和为零的分子，原独立性使每个分子为零。反过来，\(Q\) 中独立非零右理想与 \(R\) 的交都非零，因为单个非零分式可清分母；这些交仍独立。因此 \(R_R,Q_Q\) 一致维数相同。\(Q_Q\) 是有限个简单右模的直和，简单模一致，\cref{b2:prop:uniform-dimension-structure}说明其一致维数正好等于简单分量个数，也就是长度。
\end{proof}

\subsection{Goldie 假设与交换整环}
\label{b2:subsec:1902:03}

若 \(R\) 右 Noether，则全部右理想满足升链条件，特别是右零化子链稳定；\cref{b2:prop:uniform-dimension-structure}又保证右正则模一致维数有限。因此半素右 Noether 环满足 Goldie 定理。对允许非交换的整环，非零双边理想中的 \(a\ne0\) 满足 \(a^2\ne0\)，所以半素性自动成立，右 Noether 整环因而有经典右分式环；分母包含全部非零元素，使它成为除环，具体见\cref{b2:cor:domain-uniform-ore}。

\ProseParagraphBreak
交换整环甚至不必 Noether：其右正则模一致，而含非零元素的子集零化子为零，故零化子链自动稳定。Goldie 定理便恢复通常的分式域。

\ProseParagraphBreak
在 \(R=M_n(C)\)、\(C\) 为交换整环、\(n\ge1\) 的例子中，正则矩阵恰为行列式非零的矩阵。行列式非零时，伴随矩阵公式及 \(C\) 中的消去律排除两侧零化元；行列式为零时，在分式域中取非零核列，清分母得非零的 \(C\) 系数核列，将它放入一个矩阵的第一列，其余列为零，就得到非零右零化元。因此所有正则矩阵都在 \(Q=M_n(\operatorname{Frac}C)\) 中可逆，而该环的每个元素已可用非零标量矩阵作一个右分母，故它正是经典分式环。由\cref{b2:thm:artin-wedderburn}证明中的行分解，\(Q_Q=E_{11}Q\oplus\cdots\oplus E_{nn}Q\)，每个行理想都是简单右模，故 \(\ell(Q_Q)=n\)。清分母保持独立理想族的同一论证给出 \(\operatorname{udim}R_R=n\)。这里数的是 \(n\) 个简单行模，不是作为底域向量空间的 \(n^2\) 个矩阵坐标。

\ProseParagraphBreak
半素性不能删去。\cref{b2:prop:essential-without-semiprime}将对二阶上三角矩阵环给出一个本质右理想，其中没有正则元素；该环仍满足右 Goldie 条件，其经典分式环也仍是原来的非半单环。

\ProseParagraphBreak
有限一致维数也不能由另外两个条件替代。取域 \(k\)，自由代数 \(R=k\langle x,y\rangle\) 是整环，故半素；其含有非零元素的子集右零化子均为零，所以零化子升链条件成立。但右理想
\[
yR,\ xyR,\ x^2yR,\ldots
\]
独立，因为各自的基字以前缀 \(x^iy\) 开始，这些前缀彼此不会延长成另一个。故一致维数无穷，不能应用 Goldie 定理。这与上一节发现它不满足右 Ore 条件相吻合。

\begin{corollary}\label{b2:cor:domain-uniform-ore}
设 \(R\) 为非零含幺整环，不要求交换，并令 \(S=R\setminus\{0\}\)。则 \(S\) 满足右 Ore 条件，当且仅当幺性右正则模 \(R_R\) 一致。条件成立时，经典右分式环 \(RS^{-1}\) 为除环。
\end{corollary}
\begin{proof}
整环中每个非零元素都双侧正则，\(S\) 含一且乘法封闭。若 \(S\) 满足右 Ore 条件，在任意两个非零右理想 \(I,J\) 中分别取 \(a,s\ne0\)，再取 \(t\in S,b\in R\) 使 \(at=sb\)。整环性保证 \(at\ne0\)，所以 \(0\ne at\in I\cap J\)，右正则模一致。

反过来，若 \(R_R\) 一致，对 \(a,s\ne0\)，两个非零右理想 \(aR,sR\) 的交中有非零元素 \(at=sb\)。这一元素非零，保证 \(t\ne0\)，即 \(t\in S\)。若 \(a=0\)，取 \(t=1,b=0\) 即可。因此对全部 \(a\in R,s\in S\) 都满足右 Ore 条件。

由\cref{b2:thm:right-ore-localization}得到单射 \(R\hookrightarrow Q=RS^{-1}\)。每个非零分式 \(as^{-1}\) 的分子也非零，故 \(a\) 与 \(s\) 都在 \(Q\) 中可逆，其逆为 \(sa^{-1}\)。所以 \(Q\) 是除环。
\end{proof}

\begin{proposition}\label{b2:prop:essential-without-semiprime}
设 \(k\) 为域，\(R=T_2(k)\) 为含幺二阶上三角矩阵环，并令
\[
E=\left\{\begin{pmatrix}0&b\\0&c\end{pmatrix}:b,c\in k\right\}.
\]
则 \(E\) 是 \(R_R\) 的本质右理想，却不含双侧正则元素。环 \(R\) 满足右 Goldie 条件，但不半素；它的正则元素恰为单位，经典右分式环仍为 \(R\)，因而不半单。
\end{proposition}
\begin{proof}
右乘上三角矩阵保持第一列为零，所以 \(E\) 是右理想。给定非零右理想 \(I\)，取
\[
0\ne A=\begin{pmatrix}a&b\\0&c\end{pmatrix}\in I.
\]
若 \(a=0\)，则 \(A\in I\cap E\) 已非零；若 \(a\ne0\)，则 \(AE_{12}=aE_{12}\ne0\) 也属于 \(I\cap E\)。因此 \(E\) 本质。但对每个 \(A\in E\) 都有 \(AE_{11}=0\)，而 \(E_{11}\ne0\)，所以 \(E\) 没有正则元素。

环 \(R\) 作为 \(k\) 空间只有三维，故右理想升链必稳定，即 \(R\) 右 Noether；这保证右零化子升链稳定，并由\cref{b2:prop:uniform-dimension-structure}保证一致维数有限，所以 \(R\) 右 Goldie。严格上三角双边理想 \(kE_{12}\) 非零且平方为零，故 \(R\) 不半素。半单环的 Artin--Wedderburn 分解不允许非零平方为零的双边理想，因此 \(R\) 也不半单。

最后，对上面的 \(A\)，若 \(a,c\ne0\)，就有上三角逆，所以 \(A\) 正则。若 \(a=0\)，则 \(AE_{11}=0\)；若 \(c=0\)，则 \(E_{22}A=0\)。因此正则元素恰为单位。对任意 \(r\in R,s\in R^\times\)，取 \(t=1\)、\(b=s^{-1}r\)，就有 \(rt=sb\)，即满足右 Ore 条件。将已有单位求逆不增加元素，故经典右分式环就是 \(R\) 本身。
\end{proof}

% !TeX root = ../../Book2.tex
% 纲要标题：### 19.3 多项式恒等式
% 纲要位置：计划纲要.md，第 1682 行。

\section{多项式恒等式}
\label{b2:sec:1903}

上一节通过右理想和正则元寻找可以使用的分母，本节考察所有代入共同满足的乘法关系。矩阵乘法一般不交换，但固定阶数的矩阵仍满足某些共同的非交换多项式关系。多项式恒等式研究的正是这种“对任意代入都成立”的限制。把待检验的多项式写在自由代数中，就能区分一次偶然代入为零与整个代数的恒等式；外代数和 Cayley–Hamilton 定理将给出矩阵的标准恒等式。

% 纲要要点（供目录核对）：
%
% * PI 代数与标准多项式；求值同态核的交表达所有代入的共同关系，任意特征下的交替性联系外幂消失
% * 矩阵代数的 Amitsur–Levitzki 定理；Newton 生成函数明确整数可逆条件，正特征反例并入原引理；偶次外代数恢复普通行列式，整数泛型恒等式再特化至任意系数环
% * 有限维性与 PI 条件的区别；分别排除 PI 蕴含有限维或交换的误解，自由代数通过可唯一解码的字代入证明非 PI
%

\subsection{PI 代数与标准多项式}
\label{b2:subsec:1903:01}

\begin{definition}\label{b2:def:pi-algebra}
设 \(k\) 为域，\(A\) 为结合含幺 \(k\) 代数，\(m\geq1\)。对每组 \(\mathbf a=(a_1,\ldots,a_m)\in A^m\)，令
\[
\operatorname{ev}_{\mathbf a}:k\langle X_1,\ldots,X_m\rangle\longrightarrow A,
\qquad X_i\longmapsto a_i
\]
为固定 \(k\) 标量的求值同态。若非零多项式
\[
f\in k\langle X_1,\ldots,X_m\rangle
\]
对任意 \(\mathbf a\in A^m\) 都满足 \(\operatorname{ev}_{\mathbf a}(f)=f(a_1,\ldots,a_m)=0\)，即
\[
f\in\bigcap_{\mathbf a\in A^m}\ker\operatorname{ev}_{\mathbf a},
\]
则称 \(f\) 为 \(A\) 的一个\term[duoxiangshi-hengdengshi]{多项式恒等式}[polynomial identity]。若对某个 \(m\geq1\) 存在这样的非零 \(f\)，则称 \(A\) 为\term[PI-daishu]{PI 代数}[polynomial identity algebra]。
\end{definition}

这里讨论的是自由结合代数中的普通多项式恒等式，不把迹、约化迹等运算加入多项式。Cayley–Hamilton 等式中的系数随代入的矩阵而变，也不是上述意义下的一个固定多项式恒等式。

\ProseParagraphBreak
例如交换代数满足 \(X_1X_2-X_2X_1=0\)，而这个多项式在自由代数中非零。相反，“某个元素 \(a\) 满足 \(a^2=0\)”只涉及一次代入，不说明 \(X^2\) 是整个代数的恒等式。非零单位代数甚至不可能让 \(X^2\) 在所有元素上为零，因为可取 \(X=1\)。

\begin{definition}\label{b2:def:standard-polynomial}
设 \(d\geq1\) 为整数。在整数系数自由结合代数 \(\mathbb Z\langle X_1,\ldots,X_d\rangle\) 中令
\[
s_d(X_1,\ldots,X_d)
=\sum_{\sigma\in S_d}\operatorname{sgn}(\sigma)
X_{\sigma(1)}\cdots X_{\sigma(d)}.
\]
称 \(s_d\) 为 \(d\) 次\term[biaozhun-duoxiangshi]{标准多项式}[standard polynomial]。通过 \(\mathbb Z\) 到任意交换系数环的结构映射，也得到该环上的标准多项式。
\end{definition}
\symidx{\(s_d(X_1,\ldots,X_d)\)：标准多项式}

各个排列给出不同的基字，所以 \(s_d\) 在任意域上都是非零多项式。它对每个变量线性，而且交替：若两个输入相同，将每个排列与交换这两个标签后的排列配对，两个乘积相同、符号相反，因而相消。这一论证在特征二也成立，不通过除以二来判断交替性。

\begin{proposition}\label{b2:prop:finite-dimensional-pi}
设 \(k\) 为域，\(A\) 为非零结合含幺 \(k\) 代数。若 \(\dim_kA=d<\infty\)，则 \(A\) 满足标准恒等式 \(s_{d+1}=0\)。
\end{proposition}
\begin{proof}
把 \(s_{d+1}\) 的求值看成底向量空间上的交替 \((d+1)\) 线性映射 \(A^{d+1}\to A\)。由\cref{b2:prop:exterior-alternating-universal}，它通过 \(\bigwedge^{d+1}_kA\) 分解，而\cref{b2:thm:exterior-power-basis}给出 \(\bigwedge^{d+1}_kA=0\)，所以全部求值为零。交替性使重复输入的项消失；高于维数的外幂为零，正把这一计算推广到任意输入。标准多项式本身非零，故它是一个恒等式。
\end{proof}

这个维数界不是矩阵代数的最佳界。\(M_n(k)\) 的向量空间维数为 \(n^2\)，下一小节将把标准恒等式的次数从 \(n^2+1\) 降到 \(2n\)。

\subsection{Amitsur–Levitzki 定理}
\label{b2:subsec:1903:02}

标准多项式把所有排列的矩阵乘积按排列符号相加。外代数恰能把这些符号编码进乘法：在交换含幺 \(\mathbb Q\) 代数 \(B\) 上取矩阵 \(A_1,A_2\in M_n(B)\)，并在自由 \(B\) 模的外代数中取基元 \(e_1,e_2\)。矩阵 \(A_i\) 的条目属于零次部分 \(B\)，与外代数的基元交换，因而
\[
(e_1A_1+e_2A_2)^2
=e_1e_2(A_1A_2-A_2A_1).
\]
因为 \(e_i^2=0\)，重复输入的项消失；因为 \(e_2e_1=-e_1e_2\)，交换输入次序自动产生负号。对 \(2n\) 个输入组成 \(X=\sum_i e_iA_i\) 时，同样的展开使 \(X^{2n}\) 中 \(e_1\cdots e_{2n}\) 的系数恰为 \(s_{2n}(A_1,\ldots,A_{2n})\)。因此证明将先建立 \(X^{2n}=0\)，再读取这个最高外积系数。

\ProseParagraphBreak
要使 \(X^{2n}\) 为零，需要把外代数中的分次符号转为普通矩阵的特征多项式信息。矩阵 \(X^2\) 的条目属于交换的偶次部分，因而可以对它使用行列式与 Cayley–Hamilton 定理；它的各次幂迹将使特征多项式只剩最高次项。下面先在含有有理数的系数环上完成这一计算，最后用整数系数的泛型矩阵去除特征限制。这一证明方法见 Rosset 的《A new proof of the Amitsur-Levitski identity》\cite[pp.~187--188]{Rosset1976AmitsurLevitski}。

\ProseParagraphBreak
幂迹与特征多项式系数之间的关系由 Newton 恒等式表达。这里所需的只是幂迹全部为零的情形，可以用行列式生成函数 \(d(t)=\det(I-tY)\) 直接证明：它的系数给出特征多项式，而它的导数由各次幂迹决定。

\begin{lemma}\label{b2:lem:trace-zero-nilpotence}
设 \(C\) 为交换含幺 \(\mathbb Q\) 代数，\(n\geq1\) 为整数，\(Y\in M_n(C)\)。若 \(\operatorname{tr}(Y^r)=0\) 对每个整数 \(r\geq1\) 成立，则 \(\det(TI_n-Y)=T^n\)，从而 \(Y^n=0\)。另一方面，若 \(k\) 为特征 \(p>0\) 的域，则 \(I_p\in M_p(k)\) 满足全部正次幂的迹为零，却不幂零。
\end{lemma}
\begin{proof}
在 \(C[\![t]\!]\) 中，\(A(t)=I-tY\) 的逆为 \(\sum_{j\ge0}t^jY^j\)。令 \(d(t)=\det A(t)\)。对行列式的置换展开逐项求导，每次恰好微分一个条目，整理该条目的余子式，得到
\[
d'(t)=\operatorname{tr}\Bigl(\operatorname{adj}\bigl(A(t)\bigr)A'(t)\Bigr).
\]
这是多项式求导恒等式，不要求条目是数域中的数。由\cref{b2:prop:adjugate-identity}以及 \(A(t)\) 可逆，
\[
\operatorname{adj}\bigl(A(t)\bigr)=d(t)A(t)^{-1}.
\]
因此
\[
d'(t)=-d(t)\sum_{j\ge0}t^j\operatorname{tr}(Y^{j+1})=0.
\]
写
\[
d(t)=1+c_1t+\cdots+c_nt^n,
\qquad
d'(t)=c_1+2c_2t+\cdots+nc_nt^{n-1}.
\]
因此 \(d'(t)=0\) 给出 \(i c_i=0\)、\(1\le i\le n\)。这里用到 \(C\) 是 \(\mathbb Q\) 代数：每个 \(i\cdot1_C\) 可逆，所以 \(c_i=0\)，得到 \(d(t)=1\)。将 \(\det(TI-Y)\) 的系数与 \(T^n\cdot d(T^{-1})\) 比较，得到特征多项式为 \(T^n\)。对交换环 \(C\) 上的矩阵应用\cref{b2:thm:exterior-cayley-hamilton}，得 \(Y^n=0\)。

在特征 \(p>0\) 的域 \(k\) 上，取 \(Y=I_p\)，则对全部 \(r\ge1\) 都有 \(Y^r=I_p\)、\(\operatorname{tr}(Y^r)=p\cdot1_k=0\)，而 \(Y^r\ne0\)，所以它不幂零。此时
\[
d(t)=\det(I_p-tI_p)=(1-t)^p=1-t^p
\]
不是常数，导数却为零。失败的正是从 \(p c_p=0\) 消去 \(p\) 的一步。
\end{proof}

上述系数比较实际只需要 \(1,\ldots,n\) 在交换系数环中可逆；\(\mathbb Q\) 代数的假设保证了这一条件。

\begin{theorem}[Amitsur–Levitzki]\label{b2:thm:amitsur-levitzki}
设 \(R\) 为交换含幺环，\(n\geq1\) 为整数，则对任意 \(A_1,\ldots,A_{2n}\in M_n(R)\)，有
\[
s_{2n}(A_1,\ldots,A_{2n})=0.
\]
特别地，任意域上的 \(n\) 阶矩阵代数都是 PI 代数。
\end{theorem}
\begin{proof}
先令 \(B\) 为交换含幺 \(\mathbb Q\) 代数，取 \(A_1,\ldots,A_{2n}\in M_n(B)\)。在自由 \(B\) 模的外代数 \(\Lambda_B(e_1,\ldots,e_{2n})\) 上组成矩阵
\[
X=\sum_{i=1}^{2n}e_iA_i.
\]
由\cref{b2:prop:exterior-graded-commutativity}，次数 \(p,q\) 的齐次元素交换时产生符号 \((-1)^{pq}\)。因此对条目分别齐次为次数 \(p,q\) 的矩阵 \(U,V\)，有
\[
\operatorname{tr}(UV)
=\sum_{i,j}u_{ij}v_{ji}
=(-1)^{pq}\sum_{i,j}v_{ji}u_{ij}
=(-1)^{pq}\operatorname{tr}(VU).
\]
取 \(U=X\)、\(V=X^{2r-1}\)，两者次数均为奇数，得到
\(\operatorname{tr}(X^{2r})=-\operatorname{tr}(X^{2r})\)。因为二可逆，全部这些迹为零。

记 \(C=\Lambda_B^{\mathrm{even}}(e_1,\ldots,e_{2n})\)，即全部偶次分量的直和。偶数次数相加仍为偶数，零次单位也属于 \(C\)；两项偶次齐次元素交换时的符号为一，所以 \(C\) 是交换含幺 \(\mathbb Q\) 代数。\(X\) 的条目是一次元素，故 \(Y=X^2\) 的条目是二次元素，\(Y\in M_n(C)\)。这使普通行列式、特征多项式及 Cayley–Hamilton 定理可以用于 \(Y\)，而不是直接用于条目位于整个外代数中的 \(X\)。又因 \(\operatorname{tr}(Y^r)=0\) 对每个 \(r\ge1\) 成立，故可在 \(M_n(C)\) 中应用\cref{b2:lem:trace-zero-nilpotence}，得到 \(Y^n=0\)，即 \(X^{2n}=0\)。展开 \(X^{2n}\) 后，重复的 \(e_i\) 使相应项为零，其余项按排列整理；\(e_1\cdots e_{2n}\) 的矩阵系数恰为
\[
\sum_{\sigma\in S_{2n}}\operatorname{sgn}(\sigma)
A_{\sigma(1)}\cdots A_{\sigma(2n)}.
\]
外幂基定理\cref{b2:thm:exterior-power-basis}保证这个基系数必须为零。因此结论对任意交换 \(\mathbb Q\) 代数成立。

最后取 \(2n\) 个泛型整数矩阵 \(\mathcal A_i=(z_{i,rs})_{r,s=1}^n\)，其中全部条目为独立不定元，系数环记为 \(\mathbb Z[\mathbf z]\)。所得 \(s_{2n}(\mathcal A_1,\ldots,\mathcal A_{2n})\) 的每个矩阵条目都是整数系数多项式 \(F_{rs}(\mathbf z)\)。在单射
\[
\mathbb Z[\mathbf z]\hookrightarrow\mathbb Q[\mathbf z]
\]
下，上一部分证明每个 \(F_{rs}\) 的像是零多项式，因而它原来的全部整数系数就已经为零。此时对给定 \(A_i\in M_n(R)\)，使用环同态
\[
\mathbb Z[\mathbf z]\longrightarrow R,\qquad
z_{i,rs}\longmapsto (A_i)_{rs}
\]
进行特化，零整数多项式仍送到零，便得到任意特征、任意交换环上的结论。证明中除以二和正整数的步骤只发生在核验整数多项式为零的中间环 \(\mathbb Q[\mathbf z]\) 上；最终特化使用的是整数系数环到 \(R\) 的映射。
\end{proof}

这称为\term[Amitsur-Levitzki-dingli]{Amitsur–Levitzki 定理}[Amitsur–Levitzki theorem]。例如二阶矩阵满足四次标准恒等式，即二十四个有序乘积按符号求和为零；它们通常不满足二次交换恒等式。本节不讨论标准恒等式次数的最小性，只使用已经证明的 \(2n\) 次恒等式。

\subsection{有限维性与 PI 条件}
\label{b2:subsec:1903:03}

PI 性对取子代数和取商代数都保持：相同恒等式在子代数的代入只是原代入的一部分；在商代数中，先提升各输入到原代数再取商即可。因此 \(M_n(k)\) 的任意子代数、以及这些子代数的任意商，都继承至少一个非零多项式恒等式。

\ProseParagraphBreak
有限维蕴含 PI，反向不成立：多项式代数 \(k[t]\) 无限维，却满足 \(X_1X_2-X_2X_1=0\)。甚至对任意 \(d\)，\(k[t_1,\ldots,t_d]\) 都满足同一个交换恒等式，而它们的增长可以不同。PI 也不蕴含交换性：当 \(n\ge2\) 时，\(M_n(k)\) 满足刚证的标准恒等式，却有 \(E_{12}E_{21}=E_{11}\ne E_{22}=E_{21}E_{12}\)。PI 条件限制共同的乘法关系，既不强制底向量空间有限维，也不强制任意两个元素交换。

\begin{proposition}\label{b2:prop:free-algebra-not-pi}
设 \(k\) 为域，则自由结合 \(k\) 代数 \(k\langle x,y\rangle\) 不满足任何非零多项式恒等式。
\end{proposition}
\begin{proof}
为了在两个生成元中保存任意候选恒等式的全部字，需要把有限个字母 \(X_i\) 编码成 \(x,y\) 字，并使连接后的字仍能唯一解码。给定任意非零 \(f\in k\langle X_1,\ldots,X_m\rangle\)，取代入
\[
X_i\longmapsto xy^ix,\qquad 1\le i\le m.
\]
每个块 \(xy^ix\) 内部只有两个端点为 \(x\)，相邻块之间出现的 \(xx\) 正好标记分界；\(i\ge1\) 保证块内没有这种分界。因此可在每个 \(xx\) 的两个 \(x\) 之间切开，读出各块中的 \(y\) 个数，恢复原下标串。不同的 \(X_i\) 字因而被送到不同的 \(x,y\) 字，空字也只映到空字。

所以原多项式中不同基字不会在代入后合并，更不会互相抵消；至少一个非零系数仍保留，故 \(f(xyx,xy^2x,\ldots,xy^mx)\ne0\)。这为每个非零候选恒等式给出一次不为零的代入，证明结论。
\end{proof}

同一个自由代数在前面已经给出 Ore 条件的失败例子，这里又显示其不存在多项式恒等式。但两个失败的证明不同：一个比较右理想的共同倍数，另一个利用字编码保存所有多项式信息，不能从其中一项结论直接代替另一项论证。

% !TeX root = ../../Book2.tex
% 纲要标题：### 19.4 代数的增长与 Gelfand–Kirillov 维数
% 纲要位置：计划纲要.md，第 1688 行。

\section{代数的增长与 Gelfand–Kirillov 维数}
\label{b2:sec:1904}

有限生成代数也可能无限维，此时可以逐步数出短乘积产生多少个线性无关方向。生成元改变会改变每一步的具体数值，却只把“允许的乘积长度”扩大一个固定倍数。Gelfand–Kirillov 维数利用这种比较，得到不依赖有限生成组的增长指标。本节固定域 \(k\)，代数均为非零含单位元的有限生成 \(k\) 代数。

关于代数增长与 GK 维数的系统论述，参见 Krause–Lenagan 的《Growth of Algebras and Gelfand-Kirillov Dimension》\cite[Chapters~1--2]{KrauseLenagan2000GrowthAlgebras}。

% 纲要要点（供目录核对）：
%
% * 有限生成代数的增长；一允许补齐较短乘积，生成元换写只使长度扩大固定倍数，区别增长指数与向量空间维数
% * Gelfand–Kirillov 维数的基本例子；常数倍可比的幂次增长、单项式格点计数、矩阵单位比较与基本 Morita 模型，综合题区分分式环、一致维数与增长
% * 不把一般非交换环都视为可按交换方式局部化；局部化先检查有限生成性，完整说明有限分母不能生成有理函数域
% ---
%

\subsection{有限生成代数的增长}
\label{b2:subsec:1904:01}

选取有限维子空间 \(V\subseteq A\)，要求 \(1\in V\)，且 \(V\) 作为代数生成 \(A\)。令
\[
V^0=k1,\qquad
V^n=\operatorname{span}_k\{\,v_1\cdots v_n:v_i\in V\,\},
\qquad
\gamma_V(n)=\dim_kV^n.
\]
因为 \(1\in V\)，较短的乘积可以补上若干个一而写成恰有 \(n\) 项的乘积。因此 \(V^n\) 记录的是长度至多 \(n\) 的全部乘积，且
\[
V^0\subseteq V^1\subseteq\cdots,\qquad
A=\bigcup_{n\ge0}V^n,\qquad V^iV^j\subseteq V^{i+j}.
\]
每一步有限维；这组空间按乘积长度逐步覆盖整个代数。
\symidx{\(\gamma_V(n)=\dim_kV^n\)：给定生成空间的增长函数}

若改用另一组有限生成元，每个旧生成元都能写成新生成元的有限个字的线性组合。旧生成元只有有限个，因而这些字的长度有共同上界 \(c\)；一个长为 \(n\) 的旧字换写后，新字的长度至多为 \(cn\)。下面的命题把这个换写过程表述为生成空间之间的包含关系。

\begin{proposition}\label{b2:prop:growth-comparison}
设 \(k\) 为域，\(A\) 为非零有限生成含幺 \(k\) 代数，\(V,W\subseteq A\) 均为含 \(1\)、有限维且生成 \(A\) 的 \(k\) 子空间。记 \(\gamma_V(n)=\dim_kV^n\)、\(\gamma_W(n)=\dim_kW^n\)，其中 \(V^n,W^n\) 由各自空间内的 \(n\) 项乘积张成，则存在正整数 \(c,d\)，使对所有 \(n\ge1\)，
\[
\gamma_V(n)\le\gamma_W(cn),\qquad
\gamma_W(n)\le\gamma_V(dn).
\]
\end{proposition}
\begin{proof}
取 \(V\) 的有限基，每个基向量都在某个 \(W^{j_i}\) 中。因为 \(W^j\) 递增，取正整数 \(c=\max\{1,j_1,\ldots,j_r\}\)，即可使 \(V\subseteq W^c\)。于是 \(V^n\subseteq W^{cn}\)，得到第一个维数不等式。交换两者得到第二个。
\end{proof}

例如 \(A=k[t]\)，取 \(V=\operatorname{span}_k\{1,t\}\) 时，\(\gamma_V(n)=n+1\)；取 \(W=\operatorname{span}_k\{1,t,t^2\}\) 时，\(\gamma_W(n)=2n+1\)。具体函数不同，但二者都随乘积长度呈一次多项式增长，后者只相当于在前者中允许两倍长度。

\subsection{Gelfand–Kirillov 维数}
\label{b2:subsec:1904:02}

\begin{definition}\label{b2:def:gk-dimension}
设 \(k\) 为域，\(A\) 为非零有限生成含幺 \(k\) 代数。取含 \(1\)、有限维且生成 \(A\) 的 \(k\) 子空间 \(V\)，令 \(V^n\) 为 \(V\) 中 \(n\) 项乘积张成的空间，\(\gamma_V(n)=\dim_kV^n\)。定义 \(A\) 的\term[Gelfand-Kirillov-weishu]{Gelfand–Kirillov 维数}[Gelfand–Kirillov dimension]为
\[
\operatorname{GKdim}_kA
=\limsup_{n\to\infty}\frac{\log\gamma_V(n)}{\log n},
\]
该数值与 \(V\) 的选择无关；取值允许为 \(+\infty\)，简称 GK 维数。
\end{definition}
\symidx{\(\operatorname{GKdim}_kA\)：代数的 Gelfand–Kirillov 维数}

对数的底数可以任取大于一的实数，因为分子、分母在换底时乘上同一个常数。若对某个 \(d\ge0\) 和正数 \(c,C\)，充分大的 \(n\) 满足 \(cn^d\le\gamma_V(n)\le Cn^d\)，则取对数后两个界都趋于 \(d\)，故 GK 维数为 \(d\)。这里要求的是增长函数与 \(n^d\) 在常数倍范围内可比，并不要求增长函数本身是多项式。若增长有下界 \(c\lambda^n\)、\(\lambda>1\)，则同一计算给出无穷 GK 维数。

\begin{proposition}\label{b2:prop:gk-independent}
设 \(k\) 为域，\(A\) 为非零有限生成含幺 \(k\) 代数，\(V\subseteq A\) 为含 \(1\)、有限维且生成 \(A\) 的 \(k\) 子空间，\(\gamma_V(n)=\dim_kV^n\)，则定义 \(\operatorname{GKdim}_kA\) 的上极限与所选 \(V\) 无关。若存在常数 \(C>0\)、\(d\geq0\) 使 \(\gamma_V(n)\leq Cn^d\) 对所有充分大的 \(n\) 成立，则 \(\operatorname{GKdim}_kA\le d\)。
\end{proposition}
\begin{proof}
由\cref{b2:prop:growth-comparison}，对某个固定 \(c\ge1\)，
\[
\frac{\log\gamma_V(n)}{\log n}
\le
\frac{\log\gamma_W(cn)}{\log(cn)}
\cdot\frac{\log(cn)}{\log n}.
\]
第二因子趋于一，第一因子的上极限不超过在全部整数指标上取得的上极限；各项非负，故得到用 \(V\) 定义的数不超过用 \(W\) 定义的数，允许右侧为无穷。交换 \(V,W\) 得到反向不等式。若 \(\gamma_V(n)\le Cn^d\)，取对数后除以 \(\log n\)，常数项 \(\log C/\log n\) 趋于零，便得最后结论。
\end{proof}

\begin{proposition}\label{b2:prop:gk-basic-examples}
设 \(k\) 为域。每个非零有限维含幺 \(k\) 代数的 GK 维数为零；对整数 \(d\geq0\)，多项式代数 \(k[t_1,\ldots,t_d]\) 的 GK 维数为 \(d\)；对整数 \(r\geq2\)，自由结合代数 \(k\langle x_1,\ldots,x_r\rangle\) 的 GK 维数为无穷。
\end{proposition}
\begin{proof}
有限维时 \(1\le\gamma_V(n)\le\dim_kA\)，因此定义中的比值趋于零。多项式环取 \(V=\operatorname{span}_k\{1,t_1,\ldots,t_d\}\)。\(V^n\) 的基是全次数不超过 \(n\) 的单项式，它们与满足 \(a_1+\cdots+a_d\le n\) 的非负整数格点一一对应。添入剩余次数 \(a_0=n-a_1-\cdots-a_d\)，便转为计数 \(a_0+\cdots+a_d=n\) 的非负整数解，得到
\[
\gamma_V(n)=\binom{n+d}{d}.
\]
\(d>0\) 时，这关于 \(n\) 是首项为 \(n^d/d!\) 的实数值多项式，所以对数比值趋于 \(d\)；\(d=0\) 时代数为 \(k\)，回到第一种情形。

对 \(r\ge2\) 个生成元的自由代数，取由一与生成元线性生成的 \(V\)。不同字为基，故
\[
\gamma_V(n)=1+r+\cdots+r^n\ge r^n.
\]
于是 \(\log\gamma_V(n)/\log n\ge n\log r/\log n\rightarrow+\infty\)，得到无穷 GK 维数。
\end{proof}

\begin{proposition}\label{b2:prop:gk-matrix-quotient}
设 \(k\) 为域，\(A\) 为非零有限生成含幺 \(k\) 代数。对整数 \(m\geq1\)，有
\[
\operatorname{GKdim}_k M_m(A)=\operatorname{GKdim}_kA.
\]
任意含幺的有限生成 \(k\) 子代数 \(B\subseteq A\) 及任意非零商 \(k\) 代数 \(C=A/I\) 还满足
\[
\operatorname{GKdim}_kB\leq\operatorname{GKdim}_kA,\qquad
\operatorname{GKdim}_kC\leq\operatorname{GKdim}_kA.
\]
\end{proposition}
\begin{proof}
取 \(W=M_m(V)\)，它有限维、含单位矩阵，并生成矩阵代数。对 \(n\ge1\)，有 \(W^n=M_m(V^n)\)：一方向由矩阵乘法展开；反方向把每个 \(n\) 项乘积 \(v_1\cdots v_n\) 放到任意矩阵单位位置，用中间指标一将它写成 \(n\) 个 \(M_m(V)\) 矩阵的乘积。故 \(\dim_kW^n=m^2\cdot\gamma_V(n)\)。取对数时只多出常数 \(\log m^2\)，不改变上极限。

对商代数，\(V\) 的像生成它，且每次乘积空间的维数不增加，故结论成立。若子代数 \(B\) 由有限维含一空间 \(U\) 生成，取 \(c\) 使 \(U\subseteq V^c\)，则 \(\dim_kU^n\le\gamma_V(cn)\)。使用与\cref{b2:prop:gk-independent}相同的对数比较，得到子代数的不等式。
\end{proof}

\(A\) 与 \(M_m(A)\) 在第十六章给出了 Morita 等价的基本模型，这里计算出该模型中的增长指数相同；以上证明直接使用矩阵单位，并未讨论一般 Morita 等价怎样影响增长。也须区别 GK 维数与向量空间维数：\(k[t]\) 的向量空间维数无穷，GK 维数却为一；矩阵代数使每个乘积空间的维数增大固定倍数，也不改变这个增长指数。

\subsection{局部化可能改变有限生成性}
\label{b2:subsec:1904:03}

本节的 GK 维数以有限生成代数为对象，所以比较局部化前后的增长之前，须先检查局部化是否仍为有限生成代数。只加入有限个分母的逆与逆转所有非零元，在这一点上可能不同。

对 \(k[t,t^{-1}]\)，取 \(V=\operatorname{span}_k\{1,t,t^{-1}\}\)，则 \(V^n\) 以 \(t^{-n},\ldots,t^n\) 为基，增长函数为 \(2n+1\)，故 GK 维数仍为一。这个具体中心局部化保留了有限生成性，并使增长函数容易直接计算。

\ProseParagraphBreak
但把所有非零多项式都变成单位，得到的 \(k(t)\) 不是有限生成 \(k\) 代数。若有限个有理函数生成它，将其分母相乘记为非零多项式 \(d\)，所得子代数包含在 \(k[t,1/d]\) 内。多项式 \(td+1\) 非常数，且任何同时整除它与 \(d\) 的多项式也整除 \((td+1)-td=1\)，所以 \(\gcd(d,td+1)=1\)。取 \(td+1\) 的一个不可约因子 \(p\)，便有 \(p\nmid d\)。若 \(1/p\in k[t,1/d]\)，写成 \(a/d^n\) 后就有 \(d^n=pa\)，与 \(p\nmid d\) 矛盾。因此有限个分母不能产生整个分式域，不能将这里的 GK 维数定义直接用于 \(k(t)\)。

\ProseParagraphBreak
在非交换情形，还须先通过 Ore 条件才能得到本章的单分式环。自由代数的指数增长计算没有自动给出分式构造；半素右 Noether 环之所以能作经典右局部化，依据的是\cref{b2:thm:goldie-sufficiency}及其链条件论证。遇到新的代数时，生成关系、增长、PI 性及分母方程需要分别核对。下一章将给出 Weyl 代数和量子平面的具体基与增长计算，作为这些方法的进一步应用。


\begin{chaptersummary}
右 Ore 条件使分母可以通过方程 \(at=sb\) 整理到右侧。共同分母给出分式的加法，再由作用在这个加法群上的左乘算子建立结合乘法，避免未经验证地沿用交换分数法则。正则分母保证原环不被压缩；若允许零因子，单射性便可能丢失。

\ProseParagraphBreak
Goldie 理论以右理想结构保证足够多的正则元素。有限一致维数限制独立方向，零化子升链条件与半素性排除本质理想的非零左零化子；这些步骤共同推出 Ore 条件，并使分式环没有真本质右理想，最终得到半单 Artin 结构。一致维数对应分式环正则模的长度，而非原代数的基域维数。

\ProseParagraphBreak
PI 条件要求同一个非交换多项式对全部代入为零。Amitsur–Levitzki 恒等式把矩阵阶数变成这样的乘法限制，正特征结论通过整系数特化获得。GK 维数则记录短字产生的线性空间的增长，衡量另一类有限性。一致维数与增长也记录不同数据：例如 \(k[x,y]/(xy)\) 的经典分式环为 \(k(x)\times k(y)\)，一致维数为二，增长指数却为一，其具体计算见\cref{b2:ex:19-crossing-lines}。这些恒等式和增长计算都不能替代分母方程的核验；下一章将计算 Weyl 代数和量子平面的基与增长。
\end{chaptersummary}

\BookExercises

\begin{optionalexercise}\label{b2:ex:19-central-matrix-fractions}
设 \(k\) 为域。在 \(R=M_2\bigl(k[t]\bigr)\) 中，以非零标量矩阵为分母。计算
\[
E_{12}(tI)^{-1}+E_{21}\bigl((t+1)I\bigr)^{-1},\qquad
E_{12}(tI)^{-1}E_{21}\bigl((t+1)I\bigr)^{-1},
\]
并比较反序乘积。
\end{optionalexercise}
\begin{solution}
两个分母在中心，共同分母可取 \(t(t+1)I\)。和为
\[
\bigl((t+1)E_{12}+tE_{21}\bigr)\bigl(t(t+1)I\bigr)^{-1}.
\]
乘积因 \(E_{12}E_{21}=E_{11}\)，为 \(E_{11}\bigl(t(t+1)I\bigr)^{-1}\)；反序则为 \(E_{22}\bigl(t(t+1)I\bigr)^{-1}\)。分母可以交换不表示分子可以交换。由于原环嵌入分式环，两个不同矩阵单位仍不同。
\end{solution}

\begin{optionalexercise}\label{b2:ex:19-normal-element-ore}
设 \(R\) 为非零含幺环，\(s\in R\) 正则，\(\sigma:R\to R\) 为含幺环自同构，满足 \(as=s\cdot\sigma(a)\) 对每个 \(a\in R\) 成立。证明 \(S=\{s^n:n\ge0\}\) 同时满足左右 Ore 条件，并在局部化中计算 \((as^{-i})(bs^{-j})\)，其中 \(a,b\in R\)、\(i,j\ge0\) 为整数。
\end{optionalexercise}
\begin{solution}
归纳得 \(as^n=s^n\sigma^n(a)\)，因此对分子 \(a\)、分母 \(s^n\)，选择新的右分母 \(s^n\) 就满足右 Ore 方程。将 \(a\) 换成 \(\sigma^{-n}(a)\)，又得 \(s^na=\sigma^{-n}(a)s^n\)，给出左条件。正则性在取幂后保持，所以分母集满足构造定理的其余条件。

原关系在局部化中等价于 \(s^{-1}a=\sigma(a)s^{-1}\)。迭代得到 \(s^{-i}b=\sigma^i(b)s^{-i}\)，故
\[
(as^{-i})(bs^{-j})=a\sigma^i(b)s^{-(i+j)}.
\]
若忽略这里的 \(\sigma^i\)，一般会得到错误乘积。
\end{solution}

\begin{optionalexercise}\label{b2:ex:19-one-denominator-failure}
设 \(k\) 为域，\(R=k\langle x,y\rangle\)，并令
\[
B=k\langle x,y,z\rangle/(xz-1,zx-1).
\]
证明 \(B\) 满足把 \(x\) 变成单位的含幺 \(k\) 代数泛性质，且自然映射 \(R\to B\) 单射。再证明 \(B\) 中的 \(zy\) 不能写成 \(ax^{-n}\)、\(a\in R,n\ge0\)。这与\cref{b2:prop:single-element-ore-obstruction}有何关系？
\end{optionalexercise}
\begin{solution}
若含幺同态 \(f:R\to A\) 使 \(f(x)\) 可逆，则必须把 \(z\) 送到 \(f(x)^{-1}\)。这个指定满足两条关系，因而由自由代数取商的泛性质唯一给出 \(B\to A\)。

为核对单射性，把一个字按 \(y\) 的出现位置分成若干个只含 \(x,z\) 的区间。在每个区间内消去相邻的 \(xz\) 或 \(zx\)，最后只剩一个整数次幂 \(x^a\)，其中负次幂用 \(z\) 表示。得到的字唯一写成
\[
x^{a_0}yx^{a_1}y\cdots yx^{a_r},\qquad
r\ge0,\quad a_0,\ldots,a_r\in\mathbb Z.
\]
这些字确实组成商代数的基：以这样的整数列为形式基，乘法由连接两列并将接点的两个指数相加定义，再双线性延拓。三个列连接时，若中间列没有 \(y\)，两种结合次序都只把同一批整数相加；否则两端接点分别相加，因而乘法结合。该代数由 \(x,y,x^{-1}\) 生成并满足给定关系，反向也由消去过程映回 \(B\)。两映射在生成元上互逆，证明所列字线性无关。

原自由代数中的字恰好对应所有指数都非负的上述字，所以 \(R\to B\) 单射。若把 \(a\in R\) 的每个字右乘 \(x^{-n}\)，只会改变最后一个指数；每个 \(y\) 之前的指数仍非负。\(zy=x^{-1}y\) 却在第一个 \(y\) 之前有指数 \(-1\)，故由字基独立性不能得到这种表示。逆转 \(x\) 的泛性质仍可通过加生成元与关系实现，但它不能给出全部元素都是单个右分式的 Ore 构造。
\end{solution}

\begin{optionalexercise}\label{b2:ex:19-zero-divisor-localization}
设 \(R\) 为非零含幺环，\(0\ne e\in R\) 为中心幂等元。证明映射 \(R\to eR\)、\(a\mapsto ea\) 满足使 \(e\) 可逆的含幺环泛性质，并求其核，其中 \(eR\) 的单位元取为 \(e\)。再设 \(k\) 为域，证明不存在从 \(M_2(k)\) 到非零含幺环的含幺同态把 \(E_{11}\) 送成单位。
\end{optionalexercise}
\begin{solution}
中心性保证 \((ea)(eb)=eab\)，所以所给映射为满射含幺环同态。其核是 \((1-e)R\)：若 \(ea=0\)，则 \(a=(1-e)a\)，反向也显然。若含幺同态 \(f:R\to A\) 使 \(f(e)\) 为单位，幂等关系迫使 \(f(e)=1\)，故 \(f((1-e)R)=0\)。因此 \(f\) 唯一通过这个商映射分解；映射自身把 \(e\) 送到 \(eR\) 的单位元。

在矩阵情形，若 \(f(E_{11})\) 可逆，也有 \(f(E_{11})=1\)，于是 \(f(E_{22})=0\)。但
\[
E_{12}E_{22}E_{21}=E_{11},
\]
又迫使 \(f(E_{11})=0\)，与目标非零含幺矛盾。这里 \(E_{11}\) 非中心，不能把矩阵环像直积环那样仅保留一个坐标因子。
\end{solution}

\begin{optionalexercise}\label{b2:ex:19-essential-lattice}
设 \(C\) 为交换整环、\(K\) 为其分式域，\(r\ge1\) 为整数，\(N\subseteq C^r\) 为 \(C\) 子模。证明 \(N\) 本质于 \(C^r\)，当且仅当它在 \(K^r\) 中的 \(K\) 线性生成空间为全体，并求 \(\operatorname{udim}C^r\)。
\end{optionalexercise}
\begin{solution}
若 \(N\) 的 \(K\) 线性生成空间真包含于 \(K^r\)，取不在其中的非零向量，清分母后得到 \(v\in C^r\)。若 \(0\ne av\in N\)、\(a\in C\)，则 \(v=a^{-1}(av)\) 又在该线性生成空间中，矛盾。因此 \(Cv\cap N=0\)，\(N\) 不本质。

反过来，若生成空间为全体，把每个标准向量写成有限个 \(N\) 元素的 \(K\) 线性组合，再取全部系数的共同非零分母 \(d\)，得到 \(dC^r\subseteq N\)。任意非零子模 \(L\subseteq C^r\) 中取 \(0\ne x\)，则 \(dx\ne0\) 且属于 \(L\cap N\)，所以 \(N\) 本质。每个坐标模 \(C\) 一致，整个 \(C^r\) 为 \(r\) 个一致模的直和，故由\cref{b2:prop:uniform-dimension-structure}，一致维数为 \(r\)。
\end{solution}

\begin{optionalexercise}\label{b2:ex:19-uniform-versus-length}
设 \(k\) 为域，\(m,r\ge1\) 为整数，\(R=k[t]/(t^m)\)，\(M=R^r\)，并令 \(S=\bar t^{m-1}M\)。证明 \(S=\{v\in M:v\bar t=0\}\)，且右子模 \(N\subseteq M\) 本质于 \(M\) 当且仅当 \(S\subseteq N\)。据此求本质子模 \(N\) 的一致维数，并用 \(S,M\) 的长度比较说明一致维数没有记录什么信息。
\end{optionalexercise}
\begin{solution}
逐坐标展开截断多项式，乘 \(\bar t\) 为零恰好意味着每个坐标都是 \(\bar t^{m-1}\) 的标量倍，故所给两种描述相同。\(S\) 是 \(r\) 个简单右模的直和，每个 \(k\) 直线也是简单右子模。若 \(N\) 本质，它与每条这样的直线的交非零，因而包含该直线，于是包含 \(S\)。

反过来，若 \(S\subseteq N\)，任取非零子模 \(L\subseteq M\) 及 \(0\ne v\in L\)。在 \(v,v\bar t,\ldots,v\bar t^{m-1}\) 中取最后一个非零项，它被 \(\bar t\) 零化，故属于 \(L\cap S\subseteq L\cap N\)。这证明本质性。

由正文中 \(R_R\) 一致及\cref{b2:prop:uniform-dimension-structure}，\(M\) 的一致维数为 \(r\)。\(S\subseteq N\subseteq M\) 给出 \(\operatorname{udim}N=r\)。\(S\) 的长度为 \(r\)，而 \(M\) 的每一层 \(\bar t^iM/\bar t^{i+1}M\) 都是 \(r\) 个简单模的直和，共有 \(m\) 层，所以 \(\ell(M)=mr\)。当 \(m>1\) 时，两个本质子模有相同的一致维数而长度不同：独立方向内部还可以有不同数目的商层。
\end{solution}

\begin{optionalexercise}\label{b2:ex:19-matrix-annihilators}
设 \(k\) 为域，\(n\ge1\) 为整数。对 \(R=M_n(k)\) 和 \(X\subseteq R\)，用共同核 \(W=\bigcap_{A\in X}\ker A\) 描述 \(\operatorname{ann}_r(X)\)。据此检验右零化子升链条件，并计算 \(R_R\) 的一致维数和经典分式环。
\end{optionalexercise}
\begin{solution}
\(AB=0\) 对所有 \(A\in X\) 成立，当且仅当 \(B\) 的每个列都在 \(W\) 中。所以右零化子恰为这些列矩阵的集合，且可从其中允许的第一列恢复 \(W\)。严格增大的零化子链对应严格增大的子空间链，最多有 \(n\) 次严格增加，故升链条件成立。

正则右模是 \(E_{11}R,\ldots,E_{nn}R\) 的直和；每个非零行模在右矩阵作用下简单，因而一致，所以一致维数为 \(n\)。矩阵正则当且仅当可逆，奇异矩阵有非零核列并由此产生非零右零化矩阵。因此分式构造没有增加新元素，经典分式环就是 \(R\) 本身。
\end{solution}

\begin{optionalexercise}\label{b2:ex:19-product-domains}
设 \(r\ge1\) 为整数，\(C_1,\ldots,C_r\) 为交换整环，求 \(R=\prod_{i=1}^rC_i\) 的正则元素、经典分式环和一致维数。再设 \(k\) 为域，解释无限乘积 \(\prod_{i\ge1}k\) 为何不能由 Goldie 定理推出半单 Artin 分式环。
\end{optionalexercise}
\begin{solution}
一个元素正则当且仅当每个坐标都非零：若某坐标为零，相应坐标幂等元将它零化；若全部非零，逐坐标消去律排除零化元。于是分式环自然为 \(\prod_i\operatorname{Frac}(C_i)\)，有限个坐标的分式可将其分子、分母分别组成一个元素。各坐标理想一致且直和为全体，所以一致维数为 \(r\)。

无限乘积域中，任意多个非零坐标理想独立，所以一致维数无穷。正则元素恰好是所有坐标非零的元素，它们已经可逆，因此经典分式环仍为原无限乘积。令 \(I_n\) 为前 \(n\) 个坐标为零的理想，得到无限严格降链，故该环不是 Artin 环。有限一致维数条件在这里确实缺失。
\end{solution}

\begin{optionalexercise}\label{b2:ex:19-domain-uniform-ore}
设 \(R\) 为非零含幺整环，其非零元素满足右 Ore 条件，\(Q\) 为经典右分式除环。证明 \(R_R\) 本质于 \(Q_R\)，\(Q_R\) 一致，并证明左乘给出环同构 \(Q\cong\operatorname{End}_R(Q_R)\)。这里 \(Q_R\) 表示将右 \(Q\) 模 \(Q\) 限制为右 \(R\) 模。
\end{optionalexercise}
\begin{solution}
非零右 \(R\) 子模 \(L\subseteq Q\) 中取 \(0\ne q=as^{-1}\)。则 \(qs=a\ne0\) 属于 \(L\cap R\)，证明本质性。由\cref{b2:cor:domain-uniform-ore}，\(R_R\) 一致。若 \(L_1,L_2\subseteq Q_R\) 非零，则它们与 \(R\) 的交是非零右理想，故交中还有非零公共元素，所以 \(Q_R\) 也一致。

对任意右 \(R\) 线性映射 \(f:Q\to Q\)，令 \(c=f(1)\)。若 \(q=as^{-1}\)，则
\[
f(q)s=f(qs)=f(a)=ca,
\]
在 \(Q\) 中右乘 \(s^{-1}\) 得 \(f(q)=cq\)。因此每个这样的自同态都是唯一的左乘映射；反向左乘任意 \(c\in Q\) 都右 \(R\) 线性。复合左乘 \(c\) 与左乘 \(d\) 等于左乘 \(cd\)，故得到所称的环同构，无须改成反环。
\end{solution}

\begin{optionalexercise}\label{b2:ex:19-essential-without-regular}
设 \(k\) 为域，\(n\ge2\) 为整数，\(R=T_n(k)\)，\(E\subseteq R\) 为仅最后一列可能非零的矩阵所成的右理想。证明 \(E\) 本质于 \(R_R\) 且不含正则元素，计算 \(E,R_R,R/E\) 的一致维数，并说明一致维数为何不具有合成长度对短正合列的可加性。
\end{optionalexercise}
\begin{solution}
右乘上三角矩阵将 \(E\) 中的矩阵乘以其最后一个对角条目，故 \(E\) 为右理想。非零右理想 \(I\) 中取 \(A\ne0\)，选第 \(j\) 列非零，则 \(AE_{jn}\ne0\)，且仅最后一列非零。因此 \(I\cap E\ne0\)，\(E\) 本质。\(E\) 中每个矩阵都满足 \(AE_{11}=0\)，所以没有正则元素。

\(E=\bigoplus_{i=1}^n kE_{in}\) 是 \(n\) 个简单右模的直和，故一致维数为 \(n\)。它本质于 \(R_R\)，由\cref{b2:prop:uniform-dimension-structure}，\(R_R\) 的一致维数也为 \(n\)。取左上 \((n-1)\) 阶块的满射环同态 \(R\to T_{n-1}(k)\) 的核为 \(E\)，所以 \(R/E\) 的右子模就是 \(T_{n-1}(k)\) 的右理想。重复最后一列的本质子模计算，得到 \(\operatorname{udim}(R/E)=n-1\)，其中 \(n=2\) 时商为域，一致维数为一。因此在
\[
0\longrightarrow E\longrightarrow R\longrightarrow R/E\longrightarrow0
\]
中，中项一致维数为 \(n\)，两端之和却为 \(2n-1\)。本质子模已经占满原模的独立方向，商仍可能非零；这与有限长度的可加性不同。
\end{solution}

\begin{optionalexercise}\label{b2:ex:19-standard-three}
设 \(k\) 为域。在 \(M_2(k)\) 中计算 \(s_3(E_{11},E_{12},E_{21})\)，证明三次标准多项式在任意特征下都不是二阶矩阵的恒等式。四次标准多项式的情形如何？
\end{optionalexercise}
\begin{solution}
按六个排列展开，非零项分别是
\[
E_{11}E_{12}E_{21}=E_{11},\qquad
E_{12}E_{21}E_{11}=E_{11},\qquad
E_{21}E_{11}E_{12}=E_{22},
\]
这三项符号都为正，其余三项为零。因此结果为 \(2E_{11}+E_{22}\)，其右下条目为一，任意特征下都非零。由\cref{b2:thm:amitsur-levitzki}，\(s_4\) 却对所有二阶矩阵代入为零。这里具体检验的是标准多项式，不把结论扩大到任意形式的三次多项式。
\end{solution}

\begin{optionalexercise}\label{b2:ex:19-standard-recursion}
设 \(k\) 为域，\(d\ge2\) 为整数。按排列的首项分组，证明
\[
s_{d+1}(X_1,\ldots,X_{d+1})
=\sum_{i=1}^{d+1}(-1)^{i-1}X_i
s_d(X_1,\ldots,\widehat X_i,\ldots,X_{d+1}),
\]
其中帽号表示删去该变量。由此证明若 \(s_d\) 为恒等式，则所有更高次标准多项式也是恒等式。另计算 \(s_d(1,A_2,\ldots,A_d)\)。
\end{optionalexercise}
\begin{solution}
首项为 \(X_i\) 的排列，先把第 \(i\) 个变量移到最前产生符号 \((-1)^{i-1}\)，剩余变量的排列符号正好组成右端对应的 \(s_d\)。这证明公式，代入后每个求和项都含一次 \(s_d\) 的值，因此它为恒等式时，下一次也为恒等式，再归纳即可。

若把一放入第 \(j\) 个位置，乘积不变，而相对于将它放在首位的符号为 \((-1)^{j-1}\)。所以结果是 \(\sum_{j=1}^d(-1)^{j-1}\) 乘 \(s_{d-1}(A_2,\ldots,A_d)\)：\(d\) 偶数时为零，\(d\) 奇数时为 \(s_{d-1}\)。此处 \(d\ge2\)，所有配对关系在任意特征下都成立。
\end{solution}

\begin{optionalexercise}\label{b2:ex:19-no-one-variable-identity}
设 \(k\) 为无限域，\(A\ne0\) 为含单位元的 \(k\) 代数。证明 \(A\) 没有非零的一元多项式恒等式。解释这为什么不违背有限维矩阵满足 Cayley–Hamilton 定理。
\end{optionalexercise}
\begin{solution}
若 \(f(X)\in k[X]\) 在全部 \(A\) 元素上为零，特别在每个标量 \(\lambda1\) 上为零，则 \(f(\lambda)1=0\)。非零单位代数中的 \(k\) 作用忠实，所以 \(f(\lambda)=0\) 对所有 \(\lambda\in k\) 成立。无限域上的非零一元多项式只有有限多个根，因此 \(f=0\)，矛盾。Cayley–Hamilton 中使用的是每个矩阵自身的特征多项式，其系数随矩阵变化，并非同一个固定的一元多项式对全部矩阵成立。
\end{solution}

\begin{optionalexercise}\label{b2:ex:19-triangular-pi}
设 \(k\) 为域。证明 \(T_2(k)\) 满足恒等式
\[
(X_1X_2-X_2X_1)(X_3X_4-X_4X_3)=0,
\]
而 \(M_2(k)\) 不满足它。
\end{optionalexercise}
\begin{solution}
两个上三角矩阵的交换子在对角线上为零，因为对角条目在域中交换，所以它是严格上三角矩阵。两个严格上三角二阶矩阵的乘积为零，得到恒等式。对全矩阵环，令 \(X_1=X_3=E_{12}\)、\(X_2=X_4=E_{21}\)，两个交换子均为 \(E_{11}-E_{22}\)，其乘积为 \(I\ne0\)。因此子代数可以满足比全矩阵环更多的恒等式。
\end{solution}

\begin{optionalexercise}\label{b2:ex:19-trace-positive-characteristic}
设 \(k\) 为特征 \(p>0\) 的域，\(1\le n<p\)，\(Y\in M_n(k)\)。若 \(\operatorname{tr}(Y^r)=0\) 对 \(1\le r\le n\) 成立，证明 \(Y\) 幂零。解释为什么维数限制不能改为 \(n\le p\)，并在 \(p>2,n=2\) 时说明只检查 \(\operatorname{tr}Y=0\) 不够。
\end{optionalexercise}
\begin{solution}
写 \(d(t)=\det(I-tY)=1+\sum_{j=1}^n c_jt^j\)，\(c_0=1\)。迹零引理证明中的伴随矩阵与形式幂级数计算不需要除法，仍给出
\[
d'(t)=-d(t)\sum_{r\ge1}\operatorname{tr}(Y^r)t^{r-1}.
\]
比较 \(t^{j-1}\) 的系数，得到
\[
jc_j=-\sum_{r=1}^j c_{j-r}\operatorname{tr}(Y^r)=0,
\qquad 1\le j\le n.
\]
因 \(j<p\)，每个 \(j\) 在 \(k\) 中可逆，所以 \(c_j=0\)。于是 \(\chi_Y(T)=T^n\)，Cayley--Hamilton 定理给 \(Y^n=0\)。只需要前 \(n\) 个迹，因为决定 \(d(t)\) 的系数也只有这 \(n\) 个。

当 \(n=p\) 时，正文中的 \(I_p\) 反例说明最后一步不能成立：系数 \(c_p\) 不再由 \(pc_p=0\) 确定。若 \(p>2,n=2\)，矩阵 \(\operatorname{diag}(1,-1)\) 的迹为零但可逆，故只检查一次幂的迹不足以推出幂零。
\end{solution}

\begin{optionalexercise}\label{b2:ex:19-crossing-lines}
设 \(k\) 为域，令 \(R=k[x,y]/(xy)\)。计算它的经典分式环、一致维数和 GK 维数。可以先把它识别为 \(k[x]\times k[y]\) 中两坐标常数项相同的子环。
\end{optionalexercise}
\begin{solution}
商中的元素唯一写成 \(c+x\cdot p(x)+y\cdot q(y)\)，故取值于两坐标轴的映射把它识别为所述子环。一个元素正则当且仅当两个坐标都非零：二者非零时逐坐标用整环消去；若第一坐标为零，则被非零的 \((x,0)\) 零化，若第二坐标为零，则被 \((0,y)\) 零化。

所以分式环嵌入 \(k(x)\times k(y)\)。它也是满的：对 \(h(x)=f(x)/g(x)\)，其中 \(g\ne0\)，分式
\[
\frac{\bigl(x\cdot f(x),0\bigr)}{\bigl(x\cdot g(x),y\bigr)}
\]
的分子、分母都在原环中，分母正则，像为 \(\bigl(h(x),0\bigr)\)。另一坐标相同构造，再相加即可得到任意一对有理函数。因此经典分式环为 \(k(x)\times k(y)\)。

右理想 \((x)\oplus(y)\) 是本质子模：非零元素至少有一个非零坐标，乘 \(x\) 或 \(y\) 后就得到落在此和中的非零元素。两个分量分别同构于整环 \(k[x],k[y]\) 的相应一致模，所以一致维数为二。取生成空间 \(\operatorname{span}_k\{1,x,y\}\)，长度至多 \(n\) 的独立单项式为 \(1,x,\ldots,x^n,y,\ldots,y^n\)，增长为 \(2n+1\)，故 GK 维数为一。一致维数与增长指数在这个例子中不同。
\end{solution}

\begin{optionalexercise}\label{b2:ex:19-partial-laurent-growth}
设 \(k\) 为域。对 \(A=k[x,x^{-1},y]\)，取 \(V=\operatorname{span}_k\{1,x,x^{-1},y\}\)，精确计算 \(\dim_kV^n\)，并求 GK 维数。
\end{optionalexercise}
\begin{solution}
\(V^n\) 的基是 \(x^iy^j\)，其中 \(i\in\mathbb Z,j\ge0,|i|+j\le n\)。它们都可由至多 \(n\) 个指定生成元相乘得到；反向任意这样的乘积整理后也满足此界。不同 Laurent 单项式线性无关，清去一个共同的 \(x\) 幂分母即可归结为多项式单项式的无关性。因此
\[
\dim_kV^n=\sum_{j=0}^n\bigl(2(n-j)+1\bigr)=(n+1)^2,
\]
GK 维数为二。
\end{solution}

\begin{optionalexercise}\label{b2:ex:19-triangular-growth}
设 \(k\) 为域，\(d\ge0\) 为整数，\(C=k[t_1,\ldots,t_d]\)。求上三角矩阵代数 \(T_2(C)\) 的 GK 维数，说明它虽含非零幂零理想，仍可有与 \(C\) 相同的增长指数。
\end{optionalexercise}
\begin{solution}
它由有限个标量矩阵 \(t_iI\) 及矩阵单位 \(E_{11},E_{12},E_{22}\) 生成。作为 \(M_2(C)\) 的有限生成子代数，其 GK 维数不超过 \(d\)；而标量子代数 \(CI\cong C\) 又给出不小于 \(d\) 的下界，均由\cref{b2:prop:gk-matrix-quotient}得到。因此 GK 维数为 \(d\)。严格上三角矩阵理想非零且平方为零，这并不改变上述增长比较，也说明增长指数本身不排除幂零理想。
\end{solution}

\begin{optionalexercise}\label{b2:ex:19-monomial-quotient-growth}
设 \(k\) 为域。分别计算 \(A=k\langle x,y\rangle/(x^2,y^2)\) 与 \(B=k\langle x,y\rangle/(yx)\) 的 GK 维数。证明所使用的字确实组成基，而不只是一组生成元。
\end{optionalexercise}
\begin{solution}
由字关系生成的双边理想，恰是包含相应禁用连续子字的全部字的线性生成空间：左右乘字不会消除该子字，反向每个含子字的字都可分解成左字、关系、右字。因此剩余字在商中线性无关并生成。

在 \(A\) 中，剩余字必须交替出现 \(x,y\)，每个正长度恰有两个字，所以以一、两个生成元作生成空间时，增长为 \(1+2n\)，GK 维数为一。在 \(B\) 中，剩余字恰为 \(x^iy^j\)、\(i,j\ge0\)，因为出现 \(y\) 后不能再出现 \(x\)。长度至多 \(n\) 的数目为 \(\binom{n+2}{2}\)，GK 维数为二。两个代数都有两个生成元，但关系造成了不同增长。
\end{solution}

\begin{optionalexercise}\label{b2:ex:19-finite-denominator-growth}
设 \(k\) 为域，\(0\ne d(t)\in k[t]\)。证明 \(A=k[t,1/d]\) 的 GK 维数为一。须给出关于某个有限生成空间的上下界，而不是从它包含于分式域直接断言维数。
\end{optionalexercise}
\begin{solution}
取 \(V=\operatorname{span}_k\{1,t,1/d\}\)，记 \(m=\deg d\)。\(V^n\) 由 \(t^i/d^j\)、\(i+j\le n\) 线性生成。统一分母为 \(d^n\)，分子次数至多 \((m+1)n\)，所以
\[
\dim_kV^n\le(m+1)n+1.
\]
另一方面 \(1,t,\ldots,t^n\) 都在 \(V^n\) 中且线性无关，故维数至少为 \(n+1\)。取对数比值，两个界都给出极限一，因此 GK 维数为一。有限多个非零多项式作分母也可将它们相乘化为这一情形。
\end{solution}
